\chapter{动态规划、Bellman 方程、HJB 与黏性解}
\label{ch:dp-hjb}

% ════════════════════════════════════════════════════════════════
%  P4 控制部 · C1a 第 B 章（卷四《控制理论基础》）
%  源：Tutorial 01_数学/40_控制理论/30_动态规划与Bellman方程.md（离散 DP）
%       + 40_HJB方程与黏性解.md（连续 HJB / 黏性解）
%  定位：把最优控制的两条理论脊柱——离散 DP 一般理论（压缩/VI/PI/MPI/MDP）
%        与连续 HJB + 黏性解地基——讲透。LQR 当作"DP/HJB 精确可解检验"，
%        完整 Riccati 推导外引 ch:optimal-control-basics / ch:control，不重写。
%  记号归一：值函数统一 V（J 仅表代价泛函）；LQR 处 V=½xᵀPx（带 ½，与源/Bertsekas 一致）；
%        Hamiltonian 极小化 H=min_u{L+p·f}；协态 lambda=∇V=p（极小化约定）；折扣 γ=e^{-βΔt}。
%  勘误：E1(TvR!=Q-learning反例,用Baird) E2(Hopf–Lax归Hopf1965/Lax,Darbon–Osher用Hopf公式)
%        E3(Ye界log内含1/(1-γ)) E4(确等价归Joseph–Tou/Wonham非Åström)
%        E5(Doya2000=连续RL/HJB残差非Q-learning) E6(黏性解Evans1980雏形,Lions1994非单凭)
%        E7(Stratonovich需转Itô非DPP失败) E8(BRT补全Bayen) —— 全部已按正确值落地。
% ════════════════════════════════════════════════════════════════

\begin{practice}[前置自测]
开始之前，先试答下列问题。若已能流畅作答，可只速读本章的关键定理与速查卡；若卡住，正是本章要为你解决的。
\begin{enumerate}
  \item 完备度量空间上的 Banach 不动点定理怎么陈述？它为什么能一举回答“解是否存在、是否唯一、能否迭代求得”三个问题？
  \item 前面的 \cref{ch:optimal-control-basics} 用动态规划与 HJB 推出了 LQR 的 Riccati 方程。那里默认值函数光滑可微。可如果值函数在某些点\emph{不可微}（比如最短时间问题），HJB 方程 $V_t+\mathcal H(x,\nabla V)=0$ 还有意义吗？
  \item 最短时间到原点的双积分器，其最优控制是 bang--bang（只取 $\pm1$）。这样的“开关”发生在状态空间的哪条曲线上？值函数在这条曲线上还可微吗？
  \item 一维 Eikonal 方程 $|V'(x)|=1$、$V(\pm1)=0$，函数 $V(x)=1-|x|$ 与 $W(x)=|x|-1$ 都几乎处处满足它。凭什么说前者“对”、后者“错”？
  \item 强化学习里的 TD 误差 $\delta=r+\gamma V(s')-V(s)$ 与本章的 Bellman 方程是什么关系？为什么说“RL 本质上是在近似求解 Bellman 方程”？
\end{enumerate}
\end{practice}

\section*{本章目标}
学完本章，你应当能够：
\begin{enumerate}
  \item \textbf{陈述并严格证明}Bellman 最优性原理（反证 + 子策略拼接），说清它依赖的“代价可加性 + 马尔可夫性”两条命脉，并识别它失效的三类反例；
  \item \textbf{证明}Bellman 算子的 $\gamma$-压缩性，由 Banach 不动点定理推出最优值函数的存在唯一，并刻画值迭代的几何收敛率与所需迭代数；
  \item \textbf{区分}值迭代 (VI)、策略迭代 (PI) 与 Modified PI (MPI) 的收敛性质，理解“PI = Newton 法”与 Ye 强多项式界；
  \item \textbf{把}确定性 DP 嵌入 MDP 框架，写出 Bellman 期望/最优方程，并理解维数灾难为何是 MPC/iLQR/RL 三大方法的共同存在理由；
  \item \textbf{建立}从离散 Bellman 方程到连续 HJB 方程的 Taylor 极限，推导动态规划原理 (DPP) 与 HJB，并用验证定理把 HJB 的解连接到全局最优反馈；
  \item \textbf{阐明}HJB 与 PMP 的精确对偶——协态 = 值函数梯度、特征线 = Hamilton 典则方程、特征线相交 = 焦散 = 不可微；
  \item \textbf{陈述}黏性解的完整定义（测试函数从上/下触碰、超微分/次微分刻画），并\emph{亲手}验证一维 Eikonal 帐篷函数是黏性解、倒帐篷不是；
  \item \textbf{理解}比较原理与“变量翻倍”技巧、Perron 存在性、以及“值函数 = HJB 唯一黏性解”这一最优控制与 PDE 理论的会师；
  \item \textbf{掌握}单调数值格式的必要性、Barles--Souganidis 收敛定理，并知道 RL 的值函数学习就是在最小化 HJB/Bellman 残差。
\end{enumerate}

\subsection*{本章知识导航}
本章是一条“\emph{把最优控制写成一个关于值函数的方程，再问这个方程有没有解、解是什么样子}”的主线。它有两条平行的脊柱——\textbf{离散时间的动态规划}（代数、迭代、压缩）与\textbf{连续时间的 HJB}（偏微分方程、特征线、黏性解）——在“离散 $\to$ 连续极限”处缝合。结构如下：

\begin{center}
\small
\begin{tikzpicture}[
  node distance=5mm and 7mm, >=Stealth,
  blk/.style={draw, rounded corners, align=center, font=\small, inner sep=3pt, fill=main!6, draw=main!60!black},
  grp/.style={draw, rounded corners, align=center, font=\small, inner sep=3pt, fill=second!6, draw=second!60!black},
  vis/.style={draw, rounded corners, align=center, font=\small, inner sep=3pt, fill=third!8, draw=third!60!black},
  ln/.style={->, thick, gray!60!black}]
\node[blk] (op) {最优性原理\\（反证+拼接）};
\node[blk, right=of op] (bell) {Bellman 算子\\ $\gamma$-压缩};
\node[blk, right=of bell] (alg) {VI / PI / MPI\\ 收敛理论};
\node[grp, below=7mm of op] (mdp) {MDP / 随机 DP\\ 维数灾难};
\node[grp, right=of mdp] (lim) {离散$\to$连续\\ Taylor 极限};
\node[grp, right=of lim] (hjb) {DPP $\to$ HJB\\ 验证定理};
\node[vis, below=7mm of mdp] (pmp) {HJB$\leftrightarrow$PMP\\ 特征线/焦散};
\node[vis, right=of pmp] (noc) {经典解崩溃\\ 消失粘性};
\node[vis, right=of noc] (visc) {黏性解\\ 比较原理/Perron};
\draw[ln] (op)--(bell); \draw[ln] (bell)--(alg);
\draw[ln] (op.south) -- (mdp.north);
\draw[ln] (mdp)--(lim); \draw[ln] (lim)--(hjb);
\draw[ln] (hjb.south) -- ++(0,-3mm) -| (pmp.north);
\draw[ln] (pmp)--(noc); \draw[ln] (noc)--(visc);
\end{tikzpicture}
\end{center}

\noindent{}第 \ref{sec:dp-intro}--\ref{sec:dp-curse} 节是离散脊柱（任何读者的主干）；第 \ref{sec:dp-limit}--\ref{sec:dp-hjb-pmp} 节做离散到连续的过渡并推出 HJB；第 \ref{sec:dp-noclassical}--\ref{sec:dp-viscosity} 节是本章的\emph{理论高地}——黏性解；第 \ref{sec:dp-numerics} 节给数值方法的收敛地基，第 \ref{sec:dp-rl-bridge}--\ref{sec:dp-frontier} 节短接强化学习与前沿。

\subsection*{前置知识桥接与本章定位}
\cref{ch:optimal-control-basics}（最优控制基础）已用变分法、Pontryagin 极小值原理、动态规划与 HJB 四件工具,在“线性系统 + 二次代价”特例上把 LQR 的 Riccati 方程从三条路推了一遍（\cref{thm:oc-hjb,thm:oc-rde,thm:oc-pmp}）；导论 \cref{ch:control} 更早给出了 LQR 的 DP/HJB 双推导、CARE/DARE、iLQR/DDP 概念与 MPC 主定理（\cref{sec:ctrl-lqr,sec:ctrl-mpc}）。那两章把 LQR 这个“可精确求解的样板”讲透了，却都\emph{默认值函数足够光滑}，也都\emph{没有处理一般 DP 理论}（压缩映射、VI/PI 收敛、MDP）与\emph{值函数不可微}（黏性解）这两块。本章正是补全这两块：把 LQR 仅当作“DP/HJB 理论的精确可解检验”一笔带过、随用随引（\cref{sec:dp-hjb-derive}），而把笔墨投在 \cref{ch:control,ch:optimal-control-basics} 完全没有的\textbf{动态规划一般理论}与\textbf{黏性解地基}上。这与全书“导论给全景 $\to$ C0 做透基础 $\to$ A--Q 研究生密度严格化”的密度梯度一致：同一主题三层各占一档，本章站在最严格的一档。

\begin{note}
\textbf{本章记号与约定.}\;
值函数（cost-to-go）一律记 $V$：离散写 $V_k(x)$、连续写 $V(x,t)$；符号 $J$ 仅保留给“沿某给定策略的代价泛函”。\textbf{LQR 处取 $V(x)=\tfrac12 x^\top\mathbf P x$（带 $\tfrac12$）},与 Bertsekas 及本章源一致——这与导论 \cref{ch:control} 不带 $\tfrac12$ 的写法（$J^*=\mathbf x^\top\mathbf P\mathbf x$）相差一个 $\tfrac12$，引用彼处结论时须留意。Hamiltonian 取\textbf{极小化约定}
$\mathcal H(x,p)=\min_{u\in U}\{L(x,u)+p\cdot f(x,u)\}$，与 DP/HJB 自然一致；它与 \cref{ch:optimal-control-basics} 的 PMP 之间的换算为 $\mathcal H=-H_{\max}$、协态 $\lambda=-\nabla_x V=-p$（\cref{sec:dp-hjb-pmp} 详述）。\textbf{协态}（adjoint / co-state）统一译“协态”，不再切换“共态/伴随”。折扣因子离散记 $\gamma\in[0,1)$、连续记 $\beta>0$，二者满足 $\gamma=e^{-\beta\Delta t}\approx1-\beta\Delta t$。矩阵/向量随本书加粗（$\mathbf A,\mathbf B,\mathbf P$）；标量场的梯度统一写 $\nabla_x V\cdot f$（点积）、二阶 $\nabla^2 V$、时间偏导 $V_t$。本章主要在欧氏状态空间 $\mathbb R^n$ 上工作；流形状态的 DP 需测地/对数映射，见 \cref{ch:lie}，此处不展开。
\end{note}

% ════════════════════════════════════════════════════════════════
\section{两种最优性哲学：轨迹级与场级}
\label{sec:dp-intro}

\paragraph{动机：从一条轨迹到整张地图。}
\cref{ch:optimal-control-basics} 的 Pontryagin 极小值原理沿\emph{一条}最优轨迹给出一阶必要条件——协态 $\lambda^*(t)$ 只须定义在这条轨迹上。这已足够支撑 shooting、collocation、iLQR/DDP 等轨迹优化方法。但 PMP 有两个根本局限：其一，它只是\emph{必要}条件，满足 PMP 的轨迹可能是鞍点而非极小；其二，它只对\emph{单一}初始状态有效，换一个 $x_0$ 就得重解整个两点边值问题。动态规划 (DP) 走另一条路：它不求单条轨迹，而求一个标量场——\emph{值函数} $V(x)$，在每一个状态上都标出“从此处出发到终点的最小剩余代价”。一旦有了 $V$，最优控制退化为逐点查表 $u^*(x)\in\arg\min_u\{L(x,u)+V(f(x,u))\}$，给出\emph{最优反馈律}\cite{bertsekas2017dpoc1}。

\begin{remark}[一句话直觉]
PMP 像“沿一条河流找最低点”——只看到河流经过的那条路径；DP 像“画出整座山的等高线图”——对每一个位置都知道最优下山方向。前者便宜但换起点要重算，后者一次计算覆盖全局，代价是要在整个状态空间上求 $V$（这正是后文维数灾难的根源）。
\end{remark}

\paragraph{对偶全景。}
两种哲学并非对立，而是同一最优性条件的两面，其精确关系（协态 = 值函数梯度、特征线 = Hamilton 方程）将在 \cref{sec:dp-costate,sec:dp-hjb-pmp} 证明。先把全景列成对照（\cref{tab:dp-pmp-vs-dp}）。

\begin{table}[H]\centering\small
\caption{轨迹级 PMP 与场级 DP 的对偶}
\label{tab:dp-pmp-vs-dp}
\begin{tabular}{lll}
\toprule
维度 & PMP（轨迹级） & DP（场级） \\
\midrule
定义域 & 单条最优轨迹 & 整个状态空间 \\
条件类型 & 一阶必要条件 & 充分条件（全局最优） \\
核心对象 & 协态 $\lambda_k$ & 值函数 $V(x)$ \\
二者关系 & \multicolumn{2}{l}{$\lambda_k=\nabla_x V_k(x_k^*)$（\cref{sec:dp-costate}）} \\
计算代价 & $\mathcal O(n)$ ODE/差分 & $\mathcal O(N^d)$ 内存（维数灾难） \\
工程实现 & shooting、collocation & VI、PI、MPC、RL \\
适用场景 & 轨迹优化（已知初态） & 反馈控制（任意初态） \\
\bottomrule
\end{tabular}
\end{table}

这一对偶是现代机器人决策--控制--学习算法的理论脊柱：MPC 是沿退行视界反复做有限时域 DP（\cref{sec:ctrl-mpc}）；iLQR/DDP 是在名义轨迹邻域做局部二次 DP（\cref{ch:ddp-ilqr}）；强化学习全部价值类算法是从样本近似求解 Bellman 方程（\cref{sec:dp-rl-bridge}）；确定性无折扣 DP 的正向求解器即图搜索 Dijkstra/A\*（\cref{sec:plan-dijkstra}）。

\paragraph{历史脉络：从光学到强化学习。}
HJB 方程不是一次发明的，而是跨越 150 年由四代人接力。\textbf{Hamilton (1834)} 在研究几何光学与最小作用量原理时发现，粒子轨迹可由一个标量场（作用量函数 $S$）的梯度描述，建立了力学版的 Hamilton--Jacobi 方程 $S_t+H(q,\partial S/\partial q)=0$——“粒子力学 (ODE) 可由波动场论 (PDE) 等价描述”，比量子力学的波粒对偶早了近一个世纪。\textbf{Jacobi (1837)} 将其系统化：找到完全积分即可由代数关系得到力学通解。\textbf{Bellman (1957)} 在 RAND 公司研究多阶段决策时独立发现动态规划原理，其离散 Bellman 方程取连续极限即得 HJB；但 Bellman 假设值函数光滑，未处理正则性。\textbf{Crandall--Lions (1983)} 认识到大多数实际问题的值函数\emph{不可微}，于是用光滑测试函数从上/下触碰值函数、把不可微点的梯度条件“外包”给测试函数，一举解决了存在、唯一、稳定三大问题，开创非线性 PDE 理论新纪元\cite{crandall1983viscosity}。

\begin{remark}[史实小注]
通常说“Crandall--Lions 1983 创立黏性解”、“Lions 因此获 1994 Fields 奖”，这两句表述无误，但应避免绝对化。黏性解的雏形最早见于 Evans 1980，1983 年由 Crandall--Lions 系统化、1984 年由 Crandall--Evans--Lions 精化\cite{evans1984differential}；Lions 的 1994 Fields 奖引文涵盖黏性解\emph{与} Boltzmann 方程等多项工作，并非单凭黏性解。教材取标准归属即可，不必把任何一人写成“唯一发明者”。
\end{remark}

% ════════════════════════════════════════════════════════════════
\section{离散时间最优控制的标准形式}
\label{sec:dp-docp}

写 Bellman 方程之前须先界定“离散时间最优控制问题”。这不是走形式：不同的问题结构对应不同的 Bellman 方程与收敛性质。

\begin{definition}[四类标准问题]\label{def:dp-docp}
给定离散动力学，按时域与随机性分四类：
\begin{itemize}
  \item \textbf{有限时域确定性 (DOCP)}：$x_{k+1}=f_k(x_k,u_k)$，$x_k\in\mathcal X_k\subseteq\mathbb R^n$、$u_k\in\mathcal U_k(x_k)\subseteq\mathbb R^m$，求策略序列 $\pi=\{\mu_0,\dots,\mu_{N-1}\}$ 使
  \begin{equation}
    J_\pi(x_0)=\Phi(x_N)+\sum_{k=0}^{N-1}L_k\big(x_k,\mu_k(x_k)\big)
    \label{eq:dp-docp-cost}
  \end{equation}
  最小，其中 $\Phi$ 为末端代价、$L_k$ 为阶段代价。
  \item \textbf{有限时域随机 (SOCP)}：动力学含独立噪声 $x_{k+1}=f_k(x_k,u_k,w_k)$，目标取期望 $J_\pi(x_0)=\mathbb E_{w}\!\big[\Phi(x_N)+\sum_k L_k(x_k,\mu_k(x_k),w_k)\big]$。
  \item \textbf{无限时域 $\gamma$-折扣}：$N\to\infty$，代价 $\sum_{k=0}^\infty\gamma^k L(x_k,u_k,w_k)$，$\gamma\in[0,1)$，$L$ 有界。
  \item \textbf{无限时域平均代价}：目标 $\limsup_{N\to\infty}\tfrac1N\mathbb E\sum_{k=0}^{N-1}L$。
\end{itemize}
\end{definition}

倒立摆稳定对应无限时域折扣（持续运行）；抓取放置对应有限时域确定性 DOCP（有明确终止时间）；风中无人机追踪对应有限时域随机 SOCP；仓储机器人 24/7 运行对应平均代价。折扣因子 $\gamma$ 是“未来代价的时间贴现率”：在 $\gamma=0.99$、$100$\,Hz 控制下，\emph{有效视界}约 $1/(1-\gamma)=100$ 步 $=1$ 秒，恰是人类行走控制的典型 preview 时间；$\gamma=0.999$ 对应 $10$ 秒视界，适合慢速导航。$\gamma$ 既反映模型预测能力随时间衰减，也是级数收敛的数学需要。

\begin{tcolorbox}[colback=red!4,colframe=red!55!black,title=陷阱专栏：三个最常见的形式化错误]
\begin{itemize}
  \item \textbf{混淆“策略”与“控制序列”}：策略 $\mu_k(x)$ 是从状态到动作的\emph{映射}（反馈律），控制序列 $\{u_0,u_1,\dots\}$ 是开环的。DP 求的是策略。
  \item \textbf{$\gamma=1$ 时假设级数收敛}：无折扣须额外结构（吸收态、有限时域、或最短路径意义下的 proper 策略）才保代价有限——见 \cref{sec:dp-contraction} 末。
  \item \textbf{忽略动作约束 $u\in\mathcal U(x)$}：约束使 Bellman 方程的 $\min$ 变约束优化，LQR 的无约束解析解不再适用。
\end{itemize}
\end{tcolorbox}

% ════════════════════════════════════════════════════════════════
\section{Bellman 最优性原理：陈述与证明}
\label{sec:dp-optprinciple}

\paragraph{动机。}
序贯决策的最优“策略”是一个\emph{函数序列} $\{\mu_0,\mu_1,\dots\}$，搜索空间巨大。Bellman 最优性原理告诉我们\emph{不必}一次性搜索整个序列，而可逐步反向求解——这种“分而治之”的合法性正是它所保证的。

\begin{theorem}[Bellman 最优性原理]\label{thm:dp-optimality}
\emph{(Bellman 1957\cite{bellman1957dp}, p.\,83；现代形式见 Bertsekas Vol.\,I Prop.\,1.3.1\cite{bertsekas2017dpoc1}.)}
设 $\pi^*=\{\mu_0^*,\dots,\mu_{N-1}^*\}$ 是 SOCP 的最优策略。固定任意阶段 $i$ 与任意在 $\pi^*$ 下可达的状态 $x_i$，考虑以 $x_i$ 为初态、阶段 $i,\dots,N-1$ 上的子问题。则\textbf{截断策略} $\{\mu_i^*,\dots,\mu_{N-1}^*\}$ 对该子问题也最优。
\end{theorem}

Bellman 的原话是：“An optimal policy has the property that whatever the initial state and initial decision are, the remaining decisions must constitute an optimal policy with regard to the state resulting from the first decision.”

\begin{proof}[证明（反证 + 子策略拼接）]
\textbf{反设} $\{\mu_i^*,\dots,\mu_{N-1}^*\}$ 对子问题非最优，则存在子策略 $\{\mu_i',\dots,\mu_{N-1}'\}$ 使子问题成本严格更低：$J^{\mathrm{sub}}_{\pi'}(x_i)<J^{\mathrm{sub}}_{\pi^*}(x_i)$。
\textbf{构造拼接策略} $\pi'':=\{\mu_0^*,\dots,\mu_{i-1}^*,\mu_i',\dots,\mu_{N-1}'\}$，即前 $i$ 步不变、第 $i$ 步起换用更优子策略。
\textbf{分解总代价}：由代价的\emph{可加结构}与动力学的\emph{马尔可夫性}（$x_{k+1}$ 只依赖 $x_k,u_k,w_k$，不依赖更早历史），后缀代价只依赖 $x_i$、前缀代价不受后缀策略影响，故
\[
J_{\pi''}(x_0)=\underbrace{\mathbb E\Big[\sum_{k=0}^{i-1}L_k\big(x_k,\mu_k^*(x_k),w_k\big)\Big]}_{\text{前缀，与 }\pi^*\text{ 相同}}+\underbrace{\mathbb E\big[J^{\mathrm{sub}}_{\pi'}(x_i)\big]}_{<\,\mathbb E[J^{\mathrm{sub}}_{\pi^*}(x_i)]}<J_{\pi^*}(x_0),
\]
与 $\pi^*$ 最优矛盾。
\end{proof}

\begin{table}[H]\centering\small
\caption{证明所依赖的假设及其违反后果}
\label{tab:dp-op-assumptions}
\begin{tabular}{lll}
\toprule
假设 & 在证明中的作用 & 违反则 \\
\midrule
代价可加性 & 总代价可分解为前缀$+$后缀 & 乘积型 $\prod L_k$ 原理可能失效 \\
马尔可夫性 & 后缀代价只依赖 $x_i$ & 须扩充“信息状态”（POMDP） \\
$x_i$ 在 $\pi^*$ 下可达 & 拼接策略物理可实现 & 不可达状态不在讨论范围 \\
\bottomrule
\end{tabular}
\end{table}

\paragraph{三类反例与修复。}
最优性原理并非无条件成立。\textbf{乘积型代价} $J=\prod_k L_k$（$L_k>0$）：前缀的乘积因子会改变后缀的相对重要性——取对数 $\log J=\sum\log L_k$ 回到可加。\textbf{非马尔可夫}（隐含历史依赖）：决策者只能观测含噪 $y_k=h(x_k)+v_k$，最优决策依赖整个观测历史——以信念 $b_k=P(x_k\mid y_{0:k})$ 为新马尔可夫状态、在信念空间做 DP，即 POMDP 的本质。\textbf{风险敏感} $J=\mathbb E[\exp(\sum L_k)]$：指数的非线性破坏可加分解——用对数变换或 Denardo 的抽象（半环）DP。这与 CLRS \emph{Introduction to Algorithms} §15.3 的两条件“最优子结构 + 重叠子问题”是同一回事在算法教科书里的翻译，区别仅在 CLRS 面向离散组合问题、Bellman 面向连续状态控制，数学本质完全相同。

\paragraph{从原理到算法。}
由最优性原理，定义值函数 $V_k(x):=\min_{\{\mu_k,\dots\}}\mathbb E[\Phi(x_N)+\sum_{j=k}^{N-1}L_j]$（$x_k=x$），得 DP 反向递推
\begin{equation}
  V_N(x)=\Phi(x),\qquad
  V_k(x)=\min_{u\in\mathcal U_k(x)}\mathbb E_{w_k}\big[L_k(x,u,w_k)+V_{k+1}\big(f_k(x,u,w_k)\big)\big],
  \label{eq:dp-backward}
\end{equation}
最优反馈律 $\mu_k^*(x)\in\arg\min_u\{L_k(x,u)+V_{k+1}(f_k(x,u))\}$。\cref{eq:dp-backward} 是全章一切方程的母体。

% ════════════════════════════════════════════════════════════════
\section{值函数与协态的对偶}
\label{sec:dp-costate}

\paragraph{动机。}
\cref{ch:optimal-control-basics} 辛苦推导了离散 PMP 的协态方程。本节揭示一个统一两种方法的事实：\textbf{协态就是值函数的梯度}。它不仅在概念上焊接 PMP 与 DP，还解释了 iLQR 的“反向 Riccati 扫描”为何同时计算值函数的 Hessian 与梯度。

\begin{theorem}[协态 = 值函数梯度]\label{thm:dp-costate}
\emph{(Bertsekas Vol.\,I §3.3\cite{bertsekas2017dpoc1}；Kirk 2004\cite{kirk2004optimal} Ch.\,3.)}
若 $V_{k+1}$ 在 $f_k(x_k^*,u_k^*)$ 处可微、$u_k^*$ 为内点且一阶最优条件成立，则 $\lambda_k^*:=\nabla_x V_k(x_k^*)$ 满足离散 PMP 伴随方程
\begin{equation}
  \lambda_k^*=\nabla_x L_k(x_k^*,u_k^*)+\big[\nabla_x f_k(x_k^*,u_k^*)\big]^\top\lambda_{k+1}^*,
  \qquad \lambda_N^*=\nabla\Phi(x_N^*),
  \label{eq:dp-adjoint}
\end{equation}
且 $u_k^*\in\arg\min_u H_k(x_k^*,u,\lambda_{k+1}^*)$，其中 $H_k:=L_k+\lambda^\top f_k$（极小化约定）。
\end{theorem}

\begin{proof}[证明（包络定理）]
对 Bellman 方程 $V_k(x)=\min_u\{L_k(x,u)+V_{k+1}(f_k(x,u))\}$ 两侧关于 $x$ 求梯度，记取极小的控制为 $u^*(x)$：
\[
\nabla_x V_k=\nabla_x L_k+[\nabla_x f_k]^\top\nabla_x V_{k+1}
  +\underbrace{[\nabla_x u^*]^\top\big\{\nabla_u L_k+[\nabla_u f_k]^\top\nabla V_{k+1}\big\}}_{=\,0\ \text{（一阶最优条件）}}.
\]
包络项为零（因 $u^*$ 满足 $\nabla_u L_k+[\nabla_u f_k]^\top\nabla V_{k+1}=0$）。在最优轨迹处代入 $\lambda_k:=\nabla_x V_k(x_k^*)$ 即得 \cref{eq:dp-adjoint}。
\end{proof}

\begin{remark}[边际成本的物理意义]
协态 $\lambda_k$ 不是神秘的“拉格朗日乘子”：它的物理意义清晰——在状态 $x_k$ 处微扰 $\delta x_k$ 导致剩余最优代价的线性变化 $\delta V_k\approx\lambda_k^\top\delta x_k$，即“边际成本”（影子价格）。
\end{remark}

这一对偶是所有“沿轨迹反向传播梯度”算法的理论出处：iLQR/DDP 的二阶版以 $\mathbf P_k=\nabla^2 V_k$（Hessian）与 $\mathbf s_k=\nabla V_k-\mathbf P_k x_k^{\mathrm{ref}}$ 分别承载“二阶信息”与“一阶协态”（\cref{ch:ddp-ilqr}）；Neural ODE 的 adjoint 法是其连续时间版 $\lambda(t)=\nabla_x V(x^*(t),t)$。对 LQR，$V_k(x)=\tfrac12 x^\top\mathbf P_k x$ 给出 $\lambda_k=\mathbf P_k x_k^*$，恰是 \cref{ch:optimal-control-basics} 中 Riccati 解与协态的线性关系（\cref{eq:oc-lqr-u-costate}）。须强调：$\lambda_k=\nabla V_k$ \emph{只在最优轨迹上}成立，且要求 $V_k$ 在该点可微——后者在 \cref{sec:dp-noclassical} 会被打破，届时 $\lambda$ 须用次梯度替代。

% ════════════════════════════════════════════════════════════════
\section{Bellman 算子与 \texorpdfstring{$\gamma$}{gamma}-压缩}
\label{sec:dp-contraction}

\paragraph{动机。}
有限时域 DP 从 $V_N=\Phi$ 反向扫描即可。但无限时域没有“终点”，不能从无穷远反向算。问题变为：Bellman 方程 $V=\mathcal TV$ 是否有解？是否唯一？能否迭代求得？这三问的统一回答来自 Banach 不动点定理，关键在于证明 Bellman 算子是 $\gamma$-压缩的。

\begin{definition}[Bellman 算子]\label{def:dp-operator}
设状态空间 $\mathcal X$（有限或紧）、动作集 $\mathcal U(x)$ 紧、瞬时代价有界 $\|L\|_\infty\le L_{\max}$、$\gamma\in[0,1)$。值函数空间 $B(\mathcal X)=\{V:\mathcal X\to\mathbb R\mid\|V\|_\infty<\infty\}$ 在 $\sup$-范数 $\|V\|_\infty=\sup_x|V(x)|$ 下是 \textbf{Banach 空间}（完备赋范向量空间）。定义\textbf{Bellman 最优性算子} $\mathcal T$ 与\textbf{策略评估算子} $\mathcal T^\mu$：
\begin{equation}
  (\mathcal TV)(x)=\min_{u\in\mathcal U(x)}\big\{L(x,u)+\gamma\,\mathbb E_w[V(f(x,u,w))]\big\},\qquad
  (\mathcal T^\mu V)(x)=L(x,\mu(x))+\gamma\,\mathbb E_w[V(f(x,\mu(x),w))].
  \label{eq:dp-bellman-op}
\end{equation}
\end{definition}

\begin{theorem}[$\gamma$-压缩]\label{thm:dp-contraction}
\emph{(Blackwell 1965\cite{blackwell1965discounted}.)} 对任意 $V_1,V_2\in B(\mathcal X)$，
\begin{equation}
  \|\mathcal TV_1-\mathcal TV_2\|_\infty\le\gamma\,\|V_1-V_2\|_\infty,
  \label{eq:dp-contraction}
\end{equation}
$\mathcal T^\mu$ 亦然。
\end{theorem}

\begin{proof}
\textbf{引理（$\min$ 的 1-Lipschitz 性）}：对任意 $f,g:A\to\mathbb R$，$|\inf_a f-\inf_a g|\le\sup_a|f-g|$。证：取 $a_g^*$ 近似达到 $\inf_a g$，则 $\inf_a f-\inf_a g\le f(a_g^*)-g(a_g^*)\le\sup_a|f-g|$；交换 $f,g$ 得反向不等式。
\textbf{主证}：固定 $x$，令 $F_i(u)=L(x,u)+\gamma\,\mathbb E_w[V_i(f(x,u,w))]$，则 $(\mathcal TV_i)(x)=\min_u F_i(u)$。由引理 $|(\mathcal TV_1)(x)-(\mathcal TV_2)(x)|\le\sup_u|F_1(u)-F_2(u)|$。而
\[
|F_1(u)-F_2(u)|=\gamma\,\big|\mathbb E_w[V_1(f)-V_2(f)]\big|\le\gamma\,\mathbb E_w\big[|V_1(f)-V_2(f)|\big]\le\gamma\,\|V_1-V_2\|_\infty.
\]
对 $x$ 取 $\sup$ 即得 \cref{eq:dp-contraction}。
\end{proof}

\begin{corollary}[最优值函数存在唯一 + VI 几何收敛]\label{cor:dp-banach}
由 Banach 不动点定理，Bellman 最优方程 $\mathcal TV^*=V^*$ 在 $B(\mathcal X)$ 中有\emph{唯一}解 $V^*$。对任意初值 $V_0$，值迭代 $V_{k+1}=\mathcal TV_k$ 满足
\begin{equation}
  \|V_k-V^*\|_\infty\le\gamma^k\|V_0-V^*\|_\infty,
  \label{eq:dp-vi-rate}
\end{equation}
达 $\epsilon$-精度所需迭代数 $k\ge\dfrac{\log(\|V_0-V^*\|_\infty/\epsilon)}{\log(1/\gamma)}=\mathcal O\!\Big(\dfrac1{1-\gamma}\log\dfrac1\epsilon\Big)$。
\end{corollary}

\begin{example}[迭代数的量级]
$\gamma=0.99$、$\|V_0-V^*\|=100$、$\epsilon=0.01$：$k\ge\log(10^4)/\log(1/0.99)\approx9.21/0.01005\approx917$ 步。可见 $\gamma\to1$ 时 VI 极慢——这不是算法缺陷，而是问题本身的内在困难：看得越远，需越多信息才能确定最优策略。实用停止准则（Puterman Thm.\,6.3.1(c)\cite{puterman1994mdp}）：若 $\|V_{k+1}-V_k\|_\infty<(1-\gamma)\epsilon/(2\gamma)$，则贪心策略 $\mu_{k+1}$ 满足 $\|V^{\mu_{k+1}}-V^*\|_\infty<\epsilon$。
\end{example}

\begin{remark}[$\gamma=1$ 的退化]
若 $\gamma=1$，上证 Step\,2 给出 $|F_1-F_2|\le\|V_1-V_2\|_\infty$，算子仅\emph{非扩张}而非压缩。非扩张映射不保证唯一不动点、迭代可能震荡——这正是无折扣问题需额外结构（随机最短路径框架下的“proper policy”假设、或加权 $\sup$-范数 $\|V\|_w=\sup_x|V(x)|/w(x)$）的原因。
\end{remark}

% ════════════════════════════════════════════════════════════════
\section{策略迭代、Modified PI 与收敛}
\label{sec:dp-pi}

\paragraph{动机。}
VI 简单但收敛仅线性（率 $\gamma$），$\gamma=0.99$ 可能近千步。Howard 1960 的策略迭代 (PI) 在有限 MDP 上\emph{有限步终止}，实践常 5--20 步，代价是每步求解一个 $n\times n$ 线性方程组。直觉上，VI 像梯度下降（步多步廉），PI 像 Newton 法（步少步贵），MPI 介于二者之间。

\begin{definition}[Howard 策略迭代]\label{def:dp-pi}
\emph{(Howard 1960\cite{howard1960dp}.)} 从初始策略 $\mu_0$ 起，重复至 $\mu_{k+1}=\mu_k$：
\begin{enumerate}
  \item \textbf{策略评估}：精确求解 $V^{\mu_k}=\mathcal T^{\mu_k}V^{\mu_k}$。因 $\mathcal T^{\mu_k}$ 仿射，等价于线性方程组 $(\mathbf I-\gamma \mathbf P^{\mu_k})V^{\mu_k}=r^{\mu_k}$；$\gamma<1$ 使 $(\mathbf I-\gamma\mathbf P^{\mu_k})$ 严格对角占优、可逆，故 $V^{\mu_k}=(\mathbf I-\gamma\mathbf P^{\mu_k})^{-1}r^{\mu_k}$。
  \item \textbf{策略改进}：$\mu_{k+1}(x)\in\arg\min_{u}\{R(x,u)+\gamma\sum_{x'}P(x'\mid x,u)V^{\mu_k}(x')\}$（一个约束 $\arg\min$，凸性背景见 \cref{ch:nlopt}）。
\end{enumerate}
\end{definition}

\begin{theorem}[PI 有限步终止]\label{thm:dp-pi-finite}
\emph{(Puterman Thm.\,6.4.2\cite{puterman1994mdp}.)} 有限 $|\mathcal S|,|\mathcal A|$、$\gamma\in[0,1)$ 下，Howard PI 至多 $|\mathcal A|^{|\mathcal S|}$ 步终止于最优 $\mu^*$。
\end{theorem}

\begin{proof}
\textbf{改进引理}：若 $\mu_{k+1}\ne\mu_k$，则 $V^{\mu_{k+1}}\le V^{\mu_k}$ 逐点成立且至少一处严格。由改进步 $\mathcal T^{\mu_{k+1}}V^{\mu_k}\le\mathcal T^{\mu_k}V^{\mu_k}=V^{\mu_k}$；记 $W_0=V^{\mu_k}$、$W_{j+1}=\mathcal T^{\mu_{k+1}}W_j$，由 $\mathcal T^{\mu_{k+1}}$ 单调性 $W_1\le W_0$、归纳 $W_j$ 单调递减有下界，极限即不动点 $V^{\mu_{k+1}}\le V^{\mu_k}$。
\textbf{有限终止}：确定性平稳策略共 $|\mathcal A|^{|\mathcal S|}<\infty$ 个，严格改进保证不重复访问。
\textbf{终止即最优}：$\mu_{k+1}=\mu_k$ 意味着 $\mathcal TV^{\mu_k}=V^{\mu_k}$，由唯一不动点 $V^{\mu_k}=V^*$。
\end{proof}

\begin{theorem}[Ye 强多项式界 / PI = Newton]\label{thm:dp-ye-newton}
\emph{(Ye 2011\cite{ye2011simplex}.)} PI 迭代数 $\le\mathcal O\!\Big(\dfrac{mn}{1-\gamma}\log\dfrac{n}{1-\gamma}\Big)$，其中 $n=|\mathcal S|$、$m=\sum_s|\mathcal A_s|$，\emph{强多项式}（不依赖回报数值大小）。
此外（Puterman--Brumelle 1979\cite{puterman1979pi}）PI 等价于对方程 $F(V):=V-\mathcal TV=0$ 的 Newton--Kantorovich 迭代——策略评估 = 解线性化方程（Newton 方向），策略改进 = 更新线性化点，这解释了 PI 的局部二次收敛。
\end{theorem}

\begin{remark}[E3 勘误]
Ye 界的 $\log$ 内部含 $1/(1-\gamma)$——是 $\log\dfrac{n}{1-\gamma}$ 而非 $\log n$。某些速查表把它简写为 $\mathcal O(mn/(1-\gamma)\cdot\log n)$，漏掉了 $1/(1-\gamma)$，本章按完整式书写。
\end{remark}

\paragraph{Modified PI 谱系。}Puterman--Shin 1978 的 MPI 把策略评估只做 $m$ 次 $\mathcal T^{\mu_{k+1}}$ 迭代而非精确求解：$m=1$ 退化为 VI、$m=\infty$ 退化为 PI、$m$ 有限则收敛率 $\gamma^m$（比 VI 快）。这把 VI 与 PI 串成一条连续谱（\cref{tab:dp-vi-pi-mpi}），也是 Actor--Critic 的数学骨架——Critic 做若干次 TD 更新即有限次策略评估，Actor 做策略梯度即近似策略改进（\cref{sec:dp-rl-bridge}）。

\begin{table}[H]\centering\small
\caption{VI / PI / MPI 收敛性质对照}
\label{tab:dp-vi-pi-mpi}
\begin{tabular}{lcccc}
\toprule
算法 & 评估开销 & 改进开销 & 渐近收敛率 & 有限步终止 \\
\midrule
值迭代 VI & --- & $\mathcal O(mn)$ & 线性，率 $\gamma$ & 否 \\
策略迭代 PI & $\mathcal O(n^3)$ & $\mathcal O(mn)$ & 二次（局部） & 是 \\
Modified PI（$m$ 步） & $\mathcal O(mn\cdot m)$ & $\mathcal O(mn)$ & 线性，率 $\gamma^m$ & 否 \\
\bottomrule
\end{tabular}
\end{table}

\paragraph{Bellman 算子的结构性质。}$\gamma$-压缩保证收敛，但 PI 正确性与误差界还依赖另外两条结构性质，二者合起来恰是压缩性的一个充分条件。
\begin{itemize}
  \item \textbf{单调性}：$V_1\le V_2\Rightarrow\mathcal TV_1\le\mathcal TV_2$（逐点不等式两侧取 $\min_u$ 即得）。它使我们可用上下界夹逼 $V^*$：若 $V_{\mathrm{lo}}\le V^*\le V_{\mathrm{up}}$，则 $\mathcal T^kV_{\mathrm{lo}}\le V^*\le\mathcal T^kV_{\mathrm{up}}$。
  \item \textbf{平移不变性}：$\mathcal T(V+c\mathbf 1)=\mathcal TV+\gamma c\mathbf 1$。它解释了 VI 为何不需要知道 $V^*$ 的绝对值——任意常数偏移每步衰减为 $\gamma$ 倍。
  \item \textbf{Blackwell 充分条件}\cite{blackwell1965discounted}：单调性 $+$“$\gamma$-折扣” $\mathcal T(V+c\mathbf 1)\le\mathcal TV+\gamma c\mathbf 1$（$\gamma<1,c\ge0$）$\Rightarrow$ $\mathcal T$ 是 $\gamma$-压缩。这给出比直接证压缩更易验证的判据。
\end{itemize}
由不动点特征化 $V\le\mathcal TV\Rightarrow V\le V^*$（$V$ 是下界）、$V\ge\mathcal TV\Rightarrow V\ge V^*$，可导出策略次优差界（Bertsekas Vol.\,II Prop.\,2.1.1）：若 $\mu$ 对 $\tilde V$ 贪心，则 $\|V^\mu-V^*\|_\infty\le\tfrac{2\gamma}{1-\gamma}\|\tilde V-V^*\|_\infty$。对 $\gamma=0.99$ 放大因子约 $198$——这正是近似 DP/离线 RL 对值函数精度极度敏感的根源。须注意：单调性\emph{不}意味着 VI 单调收敛，$V_k$ 可能先升后降甚至震荡，但终会收敛。

% ════════════════════════════════════════════════════════════════
\section{随机 DP 与 MDP}
\label{sec:dp-mdp}

\paragraph{动机。}
把 DP 完整嵌入\textbf{马尔可夫决策过程 (MDP)}——强化学习的标准数学语言。理解 MDP 即理解 RL 的全部数学地基。

\begin{definition}[MDP 五元组]\label{def:dp-mdp}
$\mathcal M=(\mathcal S,\mathcal A,P,R,\gamma)$：状态空间 $\mathcal S$、动作空间 $\mathcal A$、转移核 $P(s'\mid s,a)$、有界奖励 $R:\mathcal S\times\mathcal A\to\mathbb R$、折扣 $\gamma\in[0,1)$。策略 $\pi:\mathcal S\to\Delta(\mathcal A)$。值函数 $V^\pi(s)=\mathbb E_\pi[\sum_{t\ge0}\gamma^tR(s_t,a_t)\mid s_0=s]$，动作值 $Q^\pi(s,a)=R(s,a)+\gamma\,\mathbb E_{s'\sim P}[V^\pi(s')]$。
\end{definition}

\textbf{Bellman 期望方程}（给定 $\pi$）与\textbf{Bellman 最优方程}：
\begin{align}
  V^\pi(s)&=\sum_a\pi(a\mid s)\Big[R(s,a)+\gamma\sum_{s'}P(s'\mid s,a)V^\pi(s')\Big],\label{eq:dp-bellman-exp}\\
  V^*(s)&=\max_a\Big\{R(s,a)+\gamma\sum_{s'}P(s'\mid s,a)V^*(s')\Big\},\qquad
  Q^*(s,a)=R(s,a)+\gamma\sum_{s'}P(s'\mid s,a)\max_{a'}Q^*(s',a').\label{eq:dp-bellman-opt}
\end{align}
确定性 DP 与 MDP 经如下字典统一：动力学 $x'=f(x,u)$ 对应 Dirac 转移核 $P(x'\mid x,u)=\delta(x'-f(x,u))$；代价最小化 $\min L$ 对应奖励最大化 $\max R=-L$（符号翻转）。\cref{sec:dp-optprinciple,sec:dp-contraction,sec:dp-pi} 的全部定理在 MDP 形式下保持不变；$Q^*$ 同样是相应压缩算子的唯一不动点。

\paragraph{风险敏感准则。}标准 MDP 优化\emph{期望}代价，但安全关键系统有时需考虑风险（\cref{tab:dp-risk}）。其中 worst-case 准则 $\min_w\sum R_t(w)$ 通向鲁棒控制与 $H_\infty$（时域 min-max Riccati 归 \cref{ch:lqr-lqg-riccati}、频域归 \cref{ch:robust-hinf}，本章不展开）。把加性噪声下的\textbf{确定性等价}（noise 仅贡献常数、不改变最优增益 $\mathbf K$，即 LQG 分离的根源）留给 \cref{ch:lqr-lqg-riccati} 的 LQG 主场；这里仅在表中点一句。

\begin{table}[H]\centering\small
\caption{代价准则：从风险中性到极端保守}
\label{tab:dp-risk}
\begin{tabular}{lll}
\toprule
准则 & 数学形式 & 适用场景 \\
\midrule
期望（风险中性） & $\mathbb E[\sum\gamma^tR_t]$ & 大多数 RL \\
CVaR（条件风险） & $\mathrm{CVaR}_\alpha[\sum R_t]$ & 安全关键系统 \\
指数效用（风险敏感） & $-\tfrac1\theta\log\mathbb E[e^{-\theta\sum R_t}]$ & 金融、鲁棒控制 \\
Worst-case（极端保守） & $\min_w\sum R_t(w)$ & $H_\infty$、博弈 \\
\bottomrule
\end{tabular}
\end{table}

\begin{remark}[E4 勘误：确定性等价的归属]
加性噪声下的确定性等价/LQG 分离最常归 Joseph--Tou 1961（确定性等价）与 Wonham 1968（分离定理严格证明），而非 Åström 1970——后者是经典教材出处而非首证，引用时不宜暗示“Åström 首次提出”。详证见 \cref{ch:lqr-lqg-riccati}。
\end{remark}

\textbf{Blackwell 最优性}\cite{blackwell1962discrete}：策略 $\pi^*$ 称 Blackwell 最优，若存在 $\gamma_0<1$ 使 $\pi^*$ 对所有 $\gamma\in[\gamma_0,1)$ 都最优；有限 MDP 存在 Blackwell 最优平稳策略。长期运行的服务机器人宜用平均代价或 Blackwell 框架，避免固定 $\gamma$ 的“折扣近视”。

% ════════════════════════════════════════════════════════════════
\section{维数灾难}
\label{sec:dp-curse}

\paragraph{动机。}
前述理论很美——Bellman 方程有唯一解、VI 必收敛、PI 有限步终止。但工程为何不直接用 tabular DP 解一切？答案是\textbf{维数灾难}。Bellman 1957 在前言首创此词，1961 年 \emph{Adaptive Control Processes} 系统展开：“It is the curse of dimensionality, a malediction that has plagued the scientist from earliest days.”

\begin{proposition}[网格化的指数爆炸]\label{prop:dp-curse}
将 $[0,1]^d$ 每维取 $N$ 格，总网格点 $|\mathcal S|=N^d$，一次完整 VI 扫描需 $\mathcal O(N^{2d}|\mathcal A|)$ 浮点运算、$\mathcal O(N^d)$ 内存。
\end{proposition}

每维 $100$ 格时：$d=2$ 为 $10^4$ 点（$80$\,KB），$d=6$ 为 $10^{12}$ 点（$8$\,TB），$d=12$ 为 $10^{24}$（不可能）。一个 6-DoF 机械臂的位形$+$速度空间 $d=12$，直接网格化完全不可行。Powell 进一步指出真实问题同时面临\emph{三重}维数灾\cite{powell2011adp}：\textbf{状态维数灾} $|\mathcal S|\sim N^{d_x}$（存 $V$ 的内存爆炸）、\textbf{动作维数灾}（连续 $\mathcal U\subseteq\mathbb R^m$ 枚举 $\arg\min_u$ 不可行，人形 $m\sim30$）、\textbf{结果维数灾}（高维噪声下算 $\mathbb E_w[\cdot]$ 的样本复杂度爆炸）。

\begin{table}[H]\centering\small
\caption{绕过维数灾的五种策略（各 1 句 + 外引）}
\label{tab:dp-mitigation}
\begin{tabular}{llll}
\toprule
策略 & 核心思想 & 代表方法 & 详见 \\
\midrule
局部二次 DP & 只在名义轨迹附近做二次逼近 & iLQR/DDP & \cref{ch:ddp-ilqr} \\
退行视界 & 只看有限步、滚动求解 & MPC & \cref{sec:ctrl-mpc} \\
函数近似 & $V\approx V_\theta$（NN/线性） & DQN、fitted VI & \cref{sec:dp-rl-bridge} \\
稀疏采样 & Monte Carlo 替穷举 & MCTS/AlphaZero & \cite{sutton2018rl} \\
路径积分 & 利用 HJB 线性化 & MPPI & \cref{ch:mppi} \\
\bottomrule
\end{tabular}
\end{table}

这三大方法家族（MPC、iLQR、RL）的\emph{存在理由}正是维数灾难——它们各以不同方式绕过它。但近似引入误差，且误差会被折扣放大：

\begin{theorem}[近似 DP 误差放大]\label{thm:dp-adp-error}
\emph{(Munos--Szepesvári 2008\cite{munos2008finite}.)} 若每步近似误差 $\|\mathcal TV_k-V_{k+1}\|_\infty\le\delta$，则 $\limsup_{k\to\infty}\|V_k-V^*\|_\infty\le\dfrac{\delta}{1-\gamma}$；对 fitted（投影 + 采样）版本，性能差被\emph{二次}放大为
\begin{equation}
  \|V^*-V^{\pi_K}\|_{p,\rho}\le\frac{2\gamma}{(1-\gamma)^2}\,C(\rho,\mu)^{1/p}\max_k\epsilon_k+\mathcal O(\gamma^K),
  \label{eq:dp-fvi-bound}
\end{equation}
其中 $C(\rho,\mu)$ 为 concentrability 系数（离线 RL 的核心理论困难）。
\end{theorem}

单步误差被 $1/(1-\gamma)$（甚至 $1/(1-\gamma)^2$）放大，对 $\gamma=0.99$ 即百倍乃至万倍——这是长视野 RL 训练困难的理论根源。

\begin{tcolorbox}[colback=red!4,colframe=red!55!black,title=陷阱专栏：“神经网络自动解维数灾”是误解（E1 勘误）]
深度网络提供表示灵活性，但统计样本复杂度仍可能随 $d$ 增长。线性函数近似下\textbf{值迭代/$Q$-learning 发散}的经典反例是 \textbf{Baird 1995}（star MDP，见 Sutton--Barto §11.2\cite{sutton2018rl}）；而 \textbf{Tsitsiklis--Van Roy 1997}\cite{tsitsiklis1997td} 的线性近似发散反例是针对 \textbf{off-policy TD(0)}（同时给出 on-policy TD 的收敛条件），\emph{不是} $Q$-learning。两者常被混为一谈，须区分：现代 RL 的成功依赖结构先验（CNN 平移不变、GNN 图结构）、海量并行采样与分布式 critic 稳定化，而非“NN 自动克服维数灾”。
\end{tcolorbox}

% ════════════════════════════════════════════════════════════════
\section{从离散 DP 到连续 HJB：极限过渡}
\label{sec:dp-limit}

\paragraph{动机。}
离散 Bellman 方程与连续时间 HJB 是同一事物的两面。理解二者的极限关系，既深化对 Bellman 方程本质的理解，又为黏性解理论铺路。这不是简单“令 $\Delta t\to0$”：连续时间带来积分代替求和、PDE 代替递推、泛函分析代替线性代数；但物理直觉不变——\emph{最优策略的尾部仍最优}。

考虑连续系统 $\dot x=f(x,u)$ 的 $\Delta t$ 离散化，离散 Bellman 方程
\[
V(x,t)=\min_u\big\{L(x,u)\Delta t+e^{-\beta\Delta t}V(x+f(x,u)\Delta t,t+\Delta t)\big\},\qquad \gamma=e^{-\beta\Delta t}.
\]
对 $V(x+f\Delta t,t+\Delta t)=V+\nabla_xV\cdot f\Delta t+V_t\Delta t+\mathcal O(\Delta t^2)$ 与 $e^{-\beta\Delta t}=1-\beta\Delta t+\mathcal O(\Delta t^2)$ 展开，代入并保留至 $\mathcal O(\Delta t)$、消去 $V$、除以 $\Delta t$、令 $\Delta t\to0$：
\begin{equation}
  \boxed{\;\beta V(x,t)=V_t+\min_u\{L(x,u)+\nabla_xV\cdot f(x,u)\}\;}
  \label{eq:dp-discount-hjb}
\end{equation}
无限时域稳态版（$V_t=0$）为 $\beta V(x)=\min_u\{L+\nabla V\cdot f\}$。这就是离散 Bellman 方程在连续极限下的形态。

\begin{table}[H]\centering\small
\caption{两种离散化范式：先离散还是先优化}
\label{tab:dp-DtO-OtD}
\begin{tabular}{llll}
\toprule
范式 & 描述 & 代表 & 理论保证 \\
\midrule
D-then-O & 先离散动力学（RK4 等）再对离散 MDP 做 DP & MPC、iLQR & 依赖离散化精度 \\
O-then-D & 先写连续 HJB PDE 再用有限差分/单调格式求解 & HJ reachability & Barles--Souganidis 收敛（\cref{sec:dp-numerics}） \\
\bottomrule
\end{tabular}
\end{table}

两范式在 $\Delta t\to0$ 时收敛到同一连续解，但有限 $\Delta t$ 下可能给出不同答案：D-then-O 的误差来自动力学离散化、O-then-D 来自 PDE 数值格式。工程中机器人本就是离散控制的，故 D-then-O 更常用。离散 Bellman 到连续 HJB 的完整对应见 \cref{tab:dp-bellman-hjb}。

\begin{table}[H]\centering\small
\caption{离散 Bellman 与连续 HJB 的逐项对应}
\label{tab:dp-bellman-hjb}
\begin{tabular}{lll}
\toprule
离散 Bellman & 连续 HJB & 对应 \\
\midrule
$V_k(x)$ & $V(x,t)$ & 值函数 \\
$\min_u\{L+V_{k+1}(f)\}$ & $\min_u\{L+\nabla V\cdot f\}$ & Bellman 算子 / Hamiltonian \\
$\gamma V_{k+1}$ & $e^{-\beta\Delta t}V\approx V-\beta V\Delta t$ & 折扣 \\
DRE（Riccati 差分） & Riccati ODE $\dot{\mathbf P}+\mathbf A^\top\mathbf P+\mathbf P\mathbf A-\mathbf P\mathbf B\mathbf R^{-1}\mathbf B^\top\mathbf P+\mathbf Q=0$ & LQR \\
$\lambda_k=\nabla V_k$ & $\lambda(t)=\nabla_xV(x^*(t),t)$ & 协态--梯度对偶（\cref{sec:dp-hjb-pmp}） \\
\bottomrule
\end{tabular}
\end{table}

% ════════════════════════════════════════════════════════════════
\section{动态规划原理与 HJB 方程推导}
\label{sec:dp-hjb-derive}

\paragraph{动机。}
\cref{sec:dp-limit} 给出了启发式极限。本节给出严格的连续时间地基：先以泛函分析的语言陈述动态规划原理 (DPP)，再由它经 Taylor 展开导出 HJB 方程。

\begin{definition}[连续 Bolza 值函数]\label{def:dp-bolza}
设 $x\in\mathbb R^n$、控制 $u\in U\subset\mathbb R^m$（$U$ 紧）、动力学 $\dot x(s)=f(x(s),u(s))$、$x(t)=x$。在容许控制类 $\mathcal U_{t,T}=\{u:[t,T]\to U\text{ 可测}\}$ 上定义
\begin{equation}
  V(x,t)=\inf_{u(\cdot)\in\mathcal U_{t,T}}\Big\{\int_t^T L(x(s),u(s))\,ds+\Phi(x(T))\Big\},\qquad V(x,T)=\Phi(x).
  \label{eq:dp-bolza}
\end{equation}
\textbf{技术假设}：$f$ 连续且对 $x$ Lipschitz（$|f(x_1,u)-f(x_2,u)|\le L_f|x_1-x_2|$，$\forall u$）、$L$ 连续有下界、$\Phi$ 连续、$U$ 紧。则对任意容许 $u$，ODE 有唯一解（Picard--Lindelöf），$V$ 在紧集上有限且连续（Bardi--Capuzzo-Dolcetta Prop.\,III.1.6\cite{bardi1997optimal}）。
\end{definition}

\begin{theorem}[动态规划原理 (DPP)]\label{thm:dp-dpp}
\emph{(Bellman 1957 连续版；Bardi--Capuzzo-Dolcetta Thm.\,III.2.1\cite{bardi1997optimal}.)} 对任意 $(x,t)$ 与 $h\in[0,T-t]$，
\begin{equation}
  V(x,t)=\inf_{u\in\mathcal U_{t,t+h}}\Big\{\int_t^{t+h}L(x(s),u(s))\,ds+V(x(t+h),t+h)\Big\}.
  \label{eq:dp-dpp}
\end{equation}
\end{theorem}

\begin{proof}[证明（两个方向）]
\textbf{方向 $\le$}：取 $[t,T]$ 上的 $\varepsilon$-最优控制 $u^\varepsilon$，则 $V(x,t)+\varepsilon\ge\int_t^{t+h}L\,ds+\int_{t+h}^T L\,ds+\Phi(x(T))$，后两项 $\ge V(x(t+h),t+h)$（$V$ 是 inf），故 $V(x,t)+\varepsilon\ge\int_t^{t+h}L\,ds+V(x(t+h),t+h)\ge$ 右端 inf，令 $\varepsilon\to0$ 得 $V\ge$ 右端。
\textbf{方向 $\ge$}：对任意 $u_1\in\mathcal U_{t,t+h}$ 与 $\varepsilon>0$，取 $u_2^\varepsilon\in\mathcal U_{t+h,T}$ 使 $V(x(t+h),t+h)+\varepsilon\ge\int_{t+h}^T L+\Phi$；拼接 $u_1,u_2^\varepsilon$ 得 $[t,T]$ 上容许控制，其总代价 $\le\int_t^{t+h}L+V(x(t+h),t+h)+\varepsilon$；对左端取 inf 得 $V(x,t)\le$ 右端 $+\varepsilon$，由 $u_1,\varepsilon$ 任意得 $V\le$ 右端。
\end{proof}

\paragraph{从 DPP 到 HJB（五步 Taylor）。}假设 $V\in C^1$（后文将看到此假设常不成立，先在光滑框架下推导）。在 \cref{eq:dp-dpp} 中取 $h$ 很小：(1) $\int_t^{t+h}L\,ds=hL(x,u)+o(h)$；(2) $V(x+hf,t+h)=V+hV_t+h\nabla_xV\cdot f+o(h)$；(3) 代入 DPP；(4) 两边消去 $V(x,t)$、除以 $h$；(5) 令 $h\to0$，注意 $V_t$ 不依赖 $u$、可提到 inf 外：
\begin{equation}
  \boxed{\;V_t(x,t)+\min_{u\in U}\{L(x,u)+\nabla_xV\cdot f(x,u)\}=0,\qquad V(x,T)=\Phi(x)\;}
  \label{eq:dp-hjb}
\end{equation}
这就是 \textbf{Hamilton--Jacobi--Bellman 方程}。定义\textbf{最优控制 Hamiltonian} $\mathcal H(x,p)=\min_{u\in U}\{L(x,u)+p\cdot f(x,u)\}$（其中 $p=\nabla_xV$ 扮演“广义动量”），HJB 紧凑写为 $V_t+\mathcal H(x,\nabla_xV)=0$。若 $\min_u$ 极值点唯一，HJB 隐含\textbf{最优反馈律} $u^*(x,t)=\arg\min_u\{L+\nabla_xV\cdot f\}$——一个状态反馈策略，比 PMP 的两点边值问题方便得多，代价是需知整个状态空间上的 $V$。

\begin{remark}[折扣稳态 HJB 的金融直觉]
当 $T=\infty$、$L,f$ 不显含时间且加折扣率 $\beta>0$，$V(x)=\inf_u\int_0^\infty e^{-\beta s}L\,ds$ 满足 $\beta V(x)=\min_u\{L+\nabla V\cdot f\}$（即 \cref{eq:dp-discount-hjb} 的稳态）。左端 $\beta V$ 是“持有值函数的利息成本”，右端是“即时代价 + 值函数改善率”——这与金融中资产定价的 no-arbitrage 条件完全同构。
\end{remark}

\begin{tcolorbox}[colback=red!4,colframe=red!55!black,title=陷阱专栏：HJB 是终值问题，且依赖 $C^1$]
\begin{itemize}
  \item \textbf{时间方向}：HJB 是\emph{终值}问题（已知 $V(x,T)=\Phi$，向 $t=0$ 反向积分），非初值问题。信息流“从未来到过去”——先知终端代价才能决定当前最优决策。数值上须反向积分，或令 $\tau=T-t$ 变为正向。
  \item \textbf{$C^1$ 假设的本质性}：第五步 Taylor 展开要求 $V\in C^1$。若 $V$ 不可微（\cref{sec:dp-noclassical} 将看到这是常态），HJB 的经典形式无意义——这正是后半章引入黏性解的核心动机，而非技术细节。
\end{itemize}
\end{tcolorbox}

\paragraph{LQR：DP/HJB 的精确可解检验。}线性 + 二次是\emph{唯一}非平凡但 HJB 可精确求解的大类：猜 $V(x,t)=\tfrac12 x^\top\mathbf P(t)x$ 代入 \cref{eq:dp-hjb}，对 $u$ 求导得 $u^*=-\mathbf R^{-1}\mathbf B^\top\mathbf P x$，回代比对系数得微分 Riccati 方程 $-\dot{\mathbf P}=\mathbf A^\top\mathbf P+\mathbf P\mathbf A-\mathbf P\mathbf B\mathbf R^{-1}\mathbf B^\top\mathbf P+\mathbf Q$，稳态即 CARE。这一推导在 \cref{ch:optimal-control-basics}（\cref{thm:oc-rde,eq:oc-care}）与导论 \cref{ch:control}（\cref{eq:ctrl-hjb,eq:ctrl-care}）已给出完整版本，本章不再重复，只把它当作“理论的可手算检验”——三条路（变分/PMP、离散 DP、HJB）殊途同归得同一 Riccati，正是 PMP、DP、HJB 三种最优性条件在 LQR 特例上的等价性。注意彼处 LQR 代价省略 $\tfrac12$（$J^*=\mathbf x^\top\mathbf P\mathbf x$），本章带 $\tfrac12$，引用时记得换算。

% ════════════════════════════════════════════════════════════════
\section{验证定理：HJB 是充分条件}
\label{sec:dp-verification}

\paragraph{动机。}
PMP 给\emph{必要}条件——满足 PMP 的轨迹未必最优（可能鞍点）。HJB 的验证定理给\emph{充分}条件：只要找到一个 $C^1$ 函数满足 HJB 与终端条件，它就是值函数、相应反馈律全局最优。这比 PMP 更强，且免除“解出后还要验证”的后顾之忧。

\begin{theorem}[验证定理]\label{thm:dp-verification}
\emph{(Fleming--Rishel 1975；Fleming--Soner Thm.\,IV.3.1\cite{flemingsoner2006}.)} 设 $W\in C^1([t_0,T]\times\mathbb R^n)$ 满足：(1) HJB 方程 $W_t+\min_u\{L+\nabla_xW\cdot f\}=0$；(2) 终端条件 $W(x,T)=\Phi(x)$；(3) 每个 $(x,t)$ 处存在 $u^*(x,t)\in\arg\min_u\{L+\nabla W\cdot f\}$；(4) 闭环 ODE $\dot x=f(x,u^*(x,t))$ 有容许解。则 $u^*$ \emph{全局最优}，且 $W\equiv V^*$。
\end{theorem}

\begin{proof}
核心思想：用 HJB 的\emph{不等式方向}证明任何容许控制不优于 $u^*$。
\textbf{Step 1}：对任意容许 $u(\cdot)$ 及轨迹 $x(\cdot)$，$\tfrac{d}{ds}W(x(s),s)=W_t+\nabla_xW\cdot f(x,u)$。
\textbf{Step 2}：由 HJB，对任意 $u\in U$ 有 $W_t+L(x,u)+\nabla W\cdot f(x,u)\ge W_t+\min_{u'}\{L+\nabla W\cdot f\}=0$，故 $\tfrac{d}{ds}W\ge-L(x(s),u(s))$。
\textbf{Step 3}：从 $t$ 到 $T$ 积分，并用 $W(x(T),T)=\Phi(x(T))$：$\int_t^T L\,ds+\Phi(x(T))\ge W(x,t)$。
\textbf{Step 4}：对所有容许 $u$ 取 inf：$V^*(x,t)\ge W(x,t)$。
\textbf{Step 5}：对 $u^*$，Step 2 的不等式取等，故 Step 3 取等：$\int_t^T L(x^*,u^*)\,ds+\Phi(x^*(T))=W(x,t)$，给出 $V^*\le W$。合并得 $V^*=W$。
\end{proof}

\begin{table}[H]\centering\small
\caption{HJB（充分）与 PMP（必要）的对偶}
\label{tab:dp-hjb-vs-pmp}
\begin{tabular}{lll}
\toprule
方面 & PMP（\cref{ch:opt-variational-pmp}） & HJB（本章） \\
\midrule
条件类型 & 必要条件 & 充分条件（经典 $C^1$ 情形） \\
对象 & 单条最优轨迹 $(x^*,\lambda^*,u^*)$ & 全状态空间值函数 $V(x,t)$ \\
维数 & $2n$ 维 ODE（两点边值） & $(n+1)$ 维 PDE \\
信息量 & 一条最优路径 & 任意初态的最优策略 \\
维数灾难 & 轻（可处理高维） & 重（$d\gtrsim6$ 经典法失效） \\
唯一性 & 多候选需筛选 & 黏性解唯一 \\
\bottomrule
\end{tabular}
\end{table}

\paragraph{从 HJB 不等式到 Lyapunov/CBF/SOS。}验证定理有一重要推论：若 $W$ 只满足\emph{不等式} $W_t+\min_u\{L+\nabla W\cdot f\}\ge0$（且 $W(x,T)\le\Phi$），则 $W\le V^*$，即 $W$ 是值函数的\emph{下界}。这把三个看似不同的框架连成一线：\textbf{Lyapunov 函数} $\dot W\le-\alpha(W)$ 相当于某种 HJB 次解条件（Lyapunov 三层定义见 \cref{sec:ctrl-lyapunov}，本章只做一句话连接）；\textbf{控制障碍函数 (CBF)} $\dot h+\alpha(h)\ge0$ 保证安全集正不变（\cref{ch:clf-cbf}）；\textbf{SOS 控制}用多项式 $W$ 满足 HJB 不等式、由 SDP 验证。须警惕两个陷阱：验证定理需 $C^1$ 正则性，而大多数 HJB 无 $C^1$ 解（\cref{sec:dp-noclassical}），届时须用黏性解版验证定理；闭环 ODE 的解须存在，当 $u^*$ 不连续（bang-bang）时闭环非经典 ODE 解，须 Filippov/Carathéodory 解。

% ════════════════════════════════════════════════════════════════
\section{HJB 与 PMP 的精确对偶}
\label{sec:dp-hjb-pmp}

\paragraph{动机。}
PMP（\cref{ch:opt-variational-pmp}）与 HJB 是最优控制的“两面”。其精确关系深刻而优美：\textbf{PMP 的协态就是值函数沿最优轨迹的梯度，Hamilton 典则方程就是 HJB 的特征线系统。}理解它，就知道 iLQR 的 backward pass 为何计算 $V_x,V_{xx}$，也知道间接法 shooting 为何等价于沿 HJB 特征线积分。

\begin{theorem}[HJB$\leftrightarrow$PMP 对偶]\label{thm:dp-duality}
设 $V\in C^2$，$x^*(\cdot),u^*(\cdot)$ 为最优。定义协态 $\lambda(t)=\nabla_xV(x^*(t),t)$。则 $\lambda$ 满足 PMP 伴随方程（极小化约定） $\dot\lambda=-L_x-f_x^\top\lambda$，且 $\lambda(t)=\nabla_xV(x^*(t),t)$。
\end{theorem}

\begin{proof}
\textbf{Step 1}：$\dot\lambda=V_{tx}+V_{xx}\dot x^*=V_{tx}+V_{xx}f(x^*,u^*)$。
\textbf{Step 2}：对 HJB $V_t+\mathcal H(x,V_x)=0$ 关于 $x$ 求偏导：$V_{tx}+\partial_x\mathcal H+\partial_p\mathcal H\cdot V_{xx}=0$，其中在 $u^*$ 处由包络定理 $\partial_p\mathcal H=f(x^*,u^*)$、$\partial_x\mathcal H=L_x+\lambda\cdot f_x$。
\textbf{Step 3}：代回 Step 1，$f$ 项相消：$\dot\lambda=-\partial_x\mathcal H=-L_x-f_x^\top\lambda$，正是 PMP 伴随方程。
\end{proof}

\paragraph{特征线视角。}一阶 PDE $V_t+\mathcal H(x,\nabla V)=0$ 的特征线 ODE 系统（经典 PDE 理论）为 $\dot x=\partial_p\mathcal H$、$\dot p=-\partial_x\mathcal H$、$\dot z=p\cdot\partial_p\mathcal H-\mathcal H$（$z=V$ 沿特征线）。前两式正是 \textbf{Hamilton 典则方程}——即 PMP 的状态--协态系统。故\textbf{PMP 是 HJB 的特征线，HJB 是 PMP 的“包络面”}，如同几何光学中费马原理（光线 = 特征线）与波动方程（波前 = PDE 等值面）的关系：光线求解容易（ODE）但只给单条路径，波前求解困难（PDE）但给整个空间的信息。

\begin{remark}[符号接口：与 A 章 PMP 的换算]
\cref{ch:opt-variational-pmp} 采用 Pontryagin 原始的\emph{极大化}约定（$H_{\max}=\lambda^\top f-L$，$u^*=\arg\max H_{\max}$），本章用\emph{极小化} $\mathcal H=\min_u\{L+p\cdot f\}$。两者是同一条件的两种等价记号：极小化 $L+\lambda\cdot f$ 等价于极大化 $\lambda^\top f-L$。换算规则一句话——\textbf{极大化协态 = 负的值函数梯度}：$\lambda_{\mathrm{PMP(max)}}=-p_{\mathrm{HJB(min)}}=-\nabla_xV$，$\mathcal H=-H_{\max}$。两章在此双向交叉引用。
\end{remark}

\paragraph{特征线相交：经典解崩溃的几何原因。}考虑从不同初态出发的最优轨迹（特征线）。当两条特征线在某点相遇，$V$ 的梯度在该点取两个不同值——$V$ 在该点\emph{不可微}。几何上，值函数等值面是特征线的“包络面”，相交处产生折叠或尖点，类似光学中的\textbf{焦散面 (caustic)}；物理上，两个不同初态经不同最优路径到达同一状态，意味着“多个最优路径”，值函数在该点连续但不可微（左右极限梯度不同）。当 $V$ 仅 Lipschitz（不可微）时，$\lambda=\nabla V$ 不唯一，对应同一状态可能有多条满足 PMP 必要条件的轨迹——PMP 不保证唯一性。这正是 \cref{sec:dp-noclassical,sec:dp-viscosity} 要深入处理的。

% ════════════════════════════════════════════════════════════════
\section{经典解为何不存在}
\label{sec:dp-noclassical}

\paragraph{动机。}
\cref{sec:dp-hjb-derive,sec:dp-verification,sec:dp-hjb-pmp} 一直假设 $V\in C^1$（甚至 $C^2$）。现在必须面对残酷现实：对绝大多数有实际意义的最优控制问题，值函数 $V$ \emph{都不是} $C^1$ 的。这不是小毛病——它使 \cref{eq:dp-hjb} 的推导严格意义上失效、验证定理无法直接用。HJB 方程 $V_t+\mathcal H(x,\nabla V)=0$ 是\emph{一阶完全非线性} PDE，与双曲守恒律（如 Burgers 方程）有相同病理：\textbf{特征线会相交}，相交处 $\nabla V$ 多值、经典解不存在。下面以四类递进的病理理解这一困难。

\paragraph{病理一：Eikonal 方程与中轴。}求到边界 $\partial\Omega$ 的距离函数：$|\nabla V|=1$（$x\in\Omega$）、$V|_{\partial\Omega}=0$，即最短时间问题（$\dot x=u$、$|u|\le1$、$L=1$）的值函数 $V(x)=\mathrm{dist}(x,\partial\Omega)$。取 $\Omega=(-1,1)\subset\mathbb R$，$V(x)=\min(x+1,1-x)$ 是“帐篷函数”，在 $x=0$ 处 $V'(0^-)=1$、$V'(0^+)=-1$，\emph{不可微}。高维时 $V$ 在\textbf{中轴}（medial axis，到 $\partial\Omega$ 最近点不唯一之处）上不可微——中轴是余维 1 集合（2D 为曲线、3D 为曲面），机器人路径规划的 Voronoi 图正是其离散版。关键观察：$-V$ 也几乎处处满足 $|\nabla(-V)|=1$，但\emph{不是}物理正确解——我们需要一种弱解概念来\emph{区分} $V$ 与 $-V$（\cref{fig:dp-tent}）。

\begin{figure}[H]
\centering
\begin{tikzpicture}[>=Stealth, scale=1.0]
  % axes
  \draw[->, gray!60!black] (-1.7,0) -- (1.9,0) node[right]{$x$};
  \draw[->, gray!60!black] (0,-1.4) -- (0,1.5) node[above]{$V$};
  \foreach \x/\lab in {-1/-1,1/1} \draw (\x,0.05)--(\x,-0.05) node[below, font=\scriptsize]{$\lab$};
  % tent (viscosity solution)
  \draw[very thick, main!70!black] (-1.5,0) -- (0,1.5) -- (1.5,0);
  \node[main!60!black, font=\small] at (1.05,1.15) {$V=\min(x{+}1,1{-}x)$};
  \node[main!60!black, font=\scriptsize, align=center] at (0,1.75) {黏性解\\（凹的一支）};
  % inverted tent (NOT viscosity solution)
  \draw[very thick, second!70!black, dashed] (-1.5,0) -- (0,-1.5) -- (1.5,0);
  \node[second!60!black, font=\small] at (1.0,-1.15) {$W=-\min(x{+}1,1{-}x)$};
  \node[second!60!black, font=\scriptsize, align=center] at (0,-1.78) {非黏性解\\（倒帐篷）};
  % test lines at apex
  \draw[third!70!black, thick] (-0.55,0.95)--(0.55,0.95);
  \node[third!60!black, font=\scriptsize] at (1.35,0.93) {$p\in D^+V=[-1,1]$};
\end{tikzpicture}
\caption{一维 Eikonal $|V'|=1$、$V(\pm1)=0$ 的两支 a.e.\ 解：帐篷 $V$（实线，凹）是黏性解，倒帐篷 $W$（虚线）不是。顶点处的超微分 $D^+V(0)=[-1,1]$ 全部满足次解条件 $|p|\le1$；谷底处 $W$ 的次微分含 $p=0$ 违反上解条件 $|p|\ge1$。博弈版（HJI/eikonal）见 \cref{sec:dg-eikonal}。}
\label{fig:dp-tent}
\end{figure}

\paragraph{病理二：双积分器 bang-bang 切换曲线。}$\ddot q=u$、$|u|\le1$、到原点 $(q,\dot q)=(0,0)$ 的最小时间。状态 $(x_1,x_2)=(q,\dot q)$：$\dot x_1=x_2$、$\dot x_2=u$。PMP 给出 bang-bang 控制 $u^*=-\mathrm{sgn}(\lambda_2(t))$，最多切换一次，切换曲线为
\begin{equation}
  \Sigma=\Big\{(x_1,x_2):x_1+\tfrac12 x_2|x_2|=0\Big\}.
  \label{eq:dp-switch-curve}
\end{equation}
最小时间值函数 $T(x_1,x_2)$ 在 $\Sigma$ 上连续但不可微：两侧特征线以不同角度汇聚到 $\Sigma$、法向导数跳跃（\cref{fig:dp-switch}）。

\begin{figure}[H]
\centering
\begin{tikzpicture}[>=Stealth, scale=1.05]
  \draw[->, gray!60!black] (-2.2,0) -- (2.3,0) node[right]{$x_1=q$};
  \draw[->, gray!60!black] (0,-1.7) -- (0,1.9) node[above]{$x_2=\dot q$};
  % switching curve x1 = -0.5 x2 |x2| : for x2>0 parabola left, x2<0 parabola right
  \draw[very thick, second!70!black, domain=0:1.7, samples=40, variable=\t]
        plot ({-0.5*\t*\t},{\t});
  \draw[very thick, second!70!black, domain=0:1.7, samples=40, variable=\t]
        plot ({0.5*\t*\t},{-\t});
  \node[second!60!black, font=\small] at (-1.55,1.55) {$\Sigma$};
  % sample parabolic trajectories u=+1 (open right) and u=-1 (open left)
  \draw[main!70!black, domain=-1.3:0.4, samples=40, variable=\s] plot ({0.5*\s*\s-0.7},{\s});
  \draw[main!70!black, domain=-0.4:1.3, samples=40, variable=\s] plot ({-0.5*\s*\s+0.7},{\s});
  \node[main!60!black, font=\scriptsize] at (1.5,1.0) {$u=-1$};
  \node[main!60!black, font=\scriptsize] at (-1.6,-0.9) {$u=+1$};
  \fill[black] (0,0) circle (1.4pt) node[below right, font=\scriptsize]{原点};
\end{tikzpicture}
\caption{双积分器最小时间问题的相平面。切换曲线 $\Sigma$（蓝，\cref{eq:dp-switch-curve}）由两段抛物线拼成；最优轨迹（红）是 $u=\pm1$ 的抛物线族，撞到 $\Sigma$ 后切换。值函数在 $\Sigma$ 上不可微。该 bang-bang 与开关曲线的完整 PMP 推导见 \cref{ch:optimal-control-basics}。}
\label{fig:dp-switch}
\end{figure}

\paragraph{病理三、四：状态约束与非光滑终端。}当最优轨迹“撞到”状态约束边界 $\partial K$（如关节限位 $q\in[-\pi,\pi]$）时，边界处“有效控制集”被约束切割，值函数在边界上不可微——需\emph{state-constrained} 黏性解（Soner 1986\cite{soner1986state}：内部用标准定义、边界只要求上解条件）。即便动力学与 running cost 都光滑，若终端代价 $\Phi$ 不光滑（如 $\Phi(x)=-|x|$），不光滑性也会沿特征线向内传播，如波前折叠。四类病理的系统总结见 \cref{tab:dp-pathology}。

\begin{table}[H]\centering\small
\caption{经典解失败的系统性原因}
\label{tab:dp-pathology}
\begin{tabular}{llll}
\toprule
原因 & 物理机制 & 典型例子 & 数学表现 \\
\midrule
特征线相交 & 多条最优路径到达同一状态 & Eikonal 中轴 & $\nabla V$ 多值 \\
Bang-bang & 最优控制不连续切换 & 时间最优 & 切换面上不可微 \\
状态约束 & 控制集被约束截断 & 关节限位 & 边界处梯度跳跃 \\
非光滑数据 & 终端/边界条件本身不光滑 & $\Phi=|x|$ & 不光滑性传播 \\
\bottomrule
\end{tabular}
\end{table}

\paragraph{消失粘性法——“黏性”之名的由来。}对有界 Lipschitz 数据，经典 $C^1$ 解一般不存在，但 $V$ 仍 Lipschitz 连续（有界梯度但不处处可微）。我们需要一种弱解：放弃处处可微、保留\emph{唯一性}（从无穷多个 a.e.\ 解中选出物理正确者）、保留 $V=V^*$。先看物理直觉：给 HJB 加人工粘性 $\varepsilon>0$，
\begin{equation}
  V_t^\varepsilon+\mathcal H(x,\nabla V^\varepsilon)=\varepsilon\,\Delta V^\varepsilon.
  \label{eq:dp-vanishing}
\end{equation}
右端 Laplacian 使方程变二阶抛物 PDE、有唯一光滑解 $V^\varepsilon$；$\varepsilon\to0$ 时 $V^\varepsilon$ 收敛到某极限 $V$——这就是\textbf{黏性解}。名称类比流体力学的粘性项：加粘性解光滑（层流），去粘性可能不光滑（激波/间断），“消失粘性”选出物理正确的弱解。Crandall--Lions 的天才在于：给出一个\emph{不依赖} $\varepsilon$ 的等价定义——纯用测试函数刻画，无须真正计算极限。

\begin{remark}[与 Burgers 方程的区别]
HJB $V_t+\mathcal H(V_x)=0$（$\mathcal H$ 凸）与 Burgers $u_t+uu_x=0$ 特征线相交时物理表现不同：Burgers 形成\emph{激波}（间断解），HJB 形成\emph{拐角}（连续但不可微解）——因为 $V$ 是 $u$ 的“原函数”，间断的 $u$ 积分成连续有拐角的 $V$。
\end{remark}

% ════════════════════════════════════════════════════════════════
\section{黏性解：严格定义与核心理论}
\label{sec:dp-viscosity}

\paragraph{动机。}
\cref{sec:dp-noclassical} 告诉我们经典解不存在。现需一个弱解概念：(1) 足够弱使解存在；(2) 足够强保证唯一；(3) 与值函数等价。Crandall--Lions 1983 的黏性解完美满足这三点。这是本章的理论高地。

\begin{definition}[黏性解]\label{def:dp-viscosity}
\emph{(Crandall--Lions 1983\cite{crandall1983viscosity}, Trans.\,AMS 277:1--42.)} 考虑一阶 PDE $F(x,V(x),\nabla V(x))=0$（对 HJB，$F(x,r,p)$ 取时间分量为 $r_t$、空间分量为 $\mathcal H(x,p_x)$）。设 $V\in C(\Omega)$：
\begin{itemize}
  \item \textbf{黏性次解}：$\forall\varphi\in C^1(\Omega)$，若 $V-\varphi$ 在 $x_0$ 取局部\emph{最大}，则 $F(x_0,V(x_0),\nabla\varphi(x_0))\le0$。
  \item \textbf{黏性上解}：$\forall\varphi\in C^1(\Omega)$，若 $V-\varphi$ 在 $x_0$ 取局部\emph{最小}，则 $F(x_0,V(x_0),\nabla\varphi(x_0))\ge0$。
  \item \textbf{黏性解}：同时为次解与上解。
\end{itemize}
\end{definition}

\paragraph{几何直觉。}在 $V$ 不可微的点 $x_0$ 无法计算 $\nabla V(x_0)$，但可找 $C^1$ 函数 $\varphi$ 从\emph{上方}或\emph{下方}触碰 $V$：从上方触碰（$V-\varphi$ 在 $x_0$ 取极大）时 $\nabla\varphi(x_0)$ 是 $V$ 的一个“超梯度”，从下方触碰时是“次梯度”。黏性解定义说的是：\emph{对所有可能的超梯度，PDE 左端 $\le0$；对所有可能的次梯度，PDE 左端 $\ge0$}。在 $V$ 可微的点，超梯度 = 次梯度 = $\nabla V$，两不等式合起来即 $F=0$，恢复经典解。这把“坏对象”（不可微函数）的性质“外包”给“好对象”（光滑测试函数），与分布理论用试探函数定义 $\delta$ 函数异曲同工，区别在分布理论用积分、黏性解用极值——反映一阶非线性 PDE 本质不同于线性 PDE。

\begin{definition}[超微分与次微分]\label{def:dp-semijet}
$D^+V(x_0)=\{p:\limsup_{y\to x_0}\frac{V(y)-V(x_0)-p\cdot(y-x_0)}{|y-x_0|}\le0\}$（超微分），$D^-V(x_0)$ 同理换 $\liminf\ge0$（次微分）。\textbf{等价刻画}：$V$ 是次解 $\iff$ $\forall x_0,\forall p\in D^+V(x_0)$，$F(x_0,V,p)\le0$；上解 $\iff$ $\forall p\in D^-V(x_0)$，$F\ge0$。当 $V$ 在 $x_0$ 可微时 $D^+V=D^-V=\{\nabla V(x_0)\}$。
\end{definition}

\begin{example}[一维 Eikonal 验证：帐篷是、倒帐篷不是]\label{ex:dp-tent}
对 $|V'|=1$、$x\in(0,1)$、$V(0)=V(1)=0$，即 $F(x,r,p)=|p|-1$。
\textbf{候选 1：}$V(x)=\min(x,1-x)$（帐篷，顶点 $x=\tfrac12$）。唯一不可微点 $x=\tfrac12$ 处 $D^+V(\tfrac12)=[-1,1]$、$D^-V(\tfrac12)=\varnothing$。次解条件：$\forall p\in[-1,1]$，$|p|-1\le0$ 成立 $\checkmark$；上解条件：空集上空真满足 $\checkmark$。$x\ne\tfrac12$ 处 $V$ 可微且 $|V'|=1$ $\checkmark$。故 $V$ 是黏性解。
\textbf{候选 2：}$W(x)=-\min(x,1-x)$（倒帐篷，谷底 $x=\tfrac12$）。$x=\tfrac12$ 处 $D^+W=\varnothing$、$D^-W=[-1,1]$。上解条件：$\forall p\in[-1,1]$ 须 $|p|-1\ge0$，但 $p=0$ 时 $|0|=0<1$，\emph{违反} $\times$。故 $W$ 不是黏性解。
\end{example}

\begin{remark}[非对称性选出物理正确解]
黏性解定义的精妙在于\emph{非对称}。对凸 Hamiltonian（Eikonal 的 $\mathcal H(p)=|p|$），它选出\emph{凹}的一支（帐篷而非倒帐篷），对应最优控制中值函数的凹性质（取 inf 保持凹性）。这不是数学上的任意约定，而是“从无穷多 a.e.\ 解中选物理正确解”的物理要求。
\end{remark}

\begin{theorem}[比较原理与唯一性]\label{thm:dp-comparison}
\emph{(Crandall--Ishii--Lions 1992\cite{crandall1992users}, Bull.\,AMS 27:1--67.)} 设 $F$ 满足“proper”条件（关于 $r$ 单调或退化椭圆）及适当结构条件：若 $u\in\mathrm{USC}(\bar\Omega)$ 为次解、$v\in\mathrm{LSC}(\bar\Omega)$ 为上解，且 $u\le v$ 于 $\partial\Omega$，则 $u\le v$ 于 $\bar\Omega$。\textbf{推论}：黏性解若存在则唯一（设 $V_1,V_2$ 皆黏性解，$V_1$ 为次解、$V_2$ 为上解 $\Rightarrow V_1\le V_2$，交换得 $V_2\le V_1$）。
\end{theorem}

\begin{proof}[证明骨架（变量翻倍）]
\textbf{Step 1（反证）}：设 $\sup_\Omega(u-v)=\delta>0$。
\textbf{Step 2（变量翻倍）}：引入 $\Phi_\alpha(x,y)=u(x)-v(y)-\tfrac\alpha2|x-y|^2$（$\alpha>0$）。由 $u$ USC、$v$ LSC、$\bar\Omega$ 紧，$\Phi_\alpha$ 在 $(\bar x_\alpha,\bar y_\alpha)$ 取最大。
\textbf{Step 3}：$\alpha\to\infty$ 时 $\alpha|\bar x_\alpha-\bar y_\alpha|^2\to0$（惩罚迫使 $\bar x_\alpha\to\bar y_\alpha$）且 $u(\bar x_\alpha)-v(\bar y_\alpha)\to\delta>0$，对大 $\alpha$ 内点。
\textbf{Step 4（构造测试函数）}：固定 $\bar y$，$x\mapsto u(x)-\tfrac\alpha2|x-\bar y|^2$ 在 $\bar x$ 取极大，次解条件给 $F(\bar x,u(\bar x),\alpha(\bar x-\bar y))\le0$；固定 $\bar x$，上解条件给 $F(\bar y,v(\bar y),\alpha(\bar x-\bar y))\ge0$。\emph{惩罚项的梯度 $\alpha(\bar x-\bar y)$ 同时充当 $u$ 与 $v$ 的测试方向}——这正是变量翻倍的妙处（解决了 $u$ 的超微分与 $v$ 的次微分在同一点方向不匹配的难题）。
\textbf{Step 5（矛盾）}：两式相减，用 $F$ 的结构条件（对 $r$ 单调或 Lipschitz），$\alpha\to\infty$ 极限得 $0\ge\delta>0$，矛盾。
\end{proof}

\begin{theorem}[Perron 存在性 + 值函数 = 唯一黏性解]\label{thm:dp-perron}
\emph{(Perron 方法；Bardi--Capuzzo-Dolcetta Thm.\,III.3.2\cite{bardi1997optimal}.)} 在适当条件下（$F$ 连续、比较原理成立、存在次解与上解），HJB 有唯一黏性解，由 $V(x)=\sup\{w(x):w\text{ 是次解},w\le\bar v\}$ 给出（$\bar v$ 为任一上解 barrier）；且最优控制问题的\emph{值函数} $V(x,t)$ 就是 HJB 的这一唯一黏性解。
\end{theorem}

\begin{proof}[证明骨架]
\textbf{Perron 存在性}：(1) $V$ 的上半连续包 $V^*$ 是次解（次解族上确界的 USC 包仍为次解，关键引理）；(2) 下半连续包 $V_*$ 是上解；(3) 由比较原理 $V^*\le V_*$；(4) 由定义 $V_*\le V^*$，故 $V^*=V_*=V$ 连续且为黏性解。
\textbf{值函数 = 黏性解}：$V$ 是次解（由 DPP \cref{eq:dp-dpp} 的 $\le$ 方向取 $\varepsilon$-最优控制、令 $h\to0$）；$V$ 是上解（由 DPP 的 $\ge$ 方向对任意控制令 $h\to0$）；唯一性由比较原理。消失粘性验证：\cref{eq:dp-vanishing} 的古典解 $V^\varepsilon$ 在 $\varepsilon\to0$ 时局部一致收敛于该唯一黏性解，证实“黏性解”名称之合理。
\end{proof}

这个定理是本章理论的会师：它把“求解最优控制问题”完全等价于“求解 HJB 的黏性解”——两者给出完全相同的函数。须留意三个陷阱：次解只要求上半连续 (USC)、上解只要求下半连续 (LSC)，黏性解二者皆满足故连续，但比较原理表述允许 $u,v$ 不连续；黏性解概念不限于一阶 PDE——Crandall--Ishii--Lions 已推广到二阶退化椭圆（随机 HJB），二阶情形用 semijets $J^{2,\pm}V$ 代替超/次微分；proper 条件需 $F$ 对 $r$ 某种单调性——纯 HJB（$F$ 不依赖 $r$）须额外强制性条件或有界域，而折扣 HJB $\beta V+\mathcal H=0$ 天然满足（$F$ 对 $r$ 严格递增）。

% ════════════════════════════════════════════════════════════════
\section{数值方法与维数灾难}
\label{sec:dp-numerics}

\paragraph{动机。}
除 LQR 等极特殊问题，HJB 通常无法解析求解，需数值方法。但 HJB 数值求解有两大挑战：一阶非线性 PDE 的特殊离散化需求（不能用标准方法），以及指数级维数灾难（$n\ge6$ 时经典网格法失效）。

\paragraph{单调性为何必要。}在笛卡尔网格上离散 $V$，对 $\nabla V$ 的差分近似\emph{必须}满足\textbf{单调性}——否则产生震荡或收敛到错误的解。这就是为何标准有限元（Galerkin/FEM）不适用：FEM 基于弱形式/变分原理，适用二阶椭圆/抛物 PDE，而 HJB 是一阶非线性 PDE，标准 Galerkin 会产生振荡（不满足单调性）。正确方法是 upwind 差分、ENO/WENO、半拉格朗日。常用的 \textbf{Lax--Friedrichs 数值 Hamiltonian} 用人工耗散保单调：
\begin{equation}
  \hat{\mathcal H}^{\mathrm{LF}}(p^+,p^-)=\mathcal H\Big(\tfrac{p^++p^-}2\Big)-\tfrac12\sum_{i=1}^n\alpha_i(p_i^+-p_i^-),
  \label{eq:dp-LF}
\end{equation}
其中 $p_i^\pm$ 为第 $i$ 方向前/后差商、$\alpha_i=\max|\partial\mathcal H/\partial p_i|$ 为耗散系数。时间离散用 TVD-RK2/RK3，受 CFL 条件 $\Delta t\le C\,\Delta x/\max_i\alpha_i$ 约束；高阶空间重构用 WENO5 平衡精度与无振荡。

\begin{theorem}[Barles--Souganidis 收敛定理]\label{thm:dp-barles}
\emph{(Barles--Souganidis 1991\cite{barles1991convergence}.)} 数值格式 $S^\rho(x,V^\rho,V^\rho)=0$（$\rho$ 为网格参数）若同时满足：\textbf{单调性}（增大邻居值不减小方程左端）、\textbf{稳定性}（$\|V^\rho\|_\infty\le C$ 一致有界）、\textbf{一致性}（对光滑 $\varphi$，$S^\rho(x,\varphi,\varphi)\to F(x,\varphi,\nabla\varphi)$），且极限方程具强比较原理，则 $V^\rho$ 局部一致收敛于唯一黏性解。
\end{theorem}

\begin{remark}[“金科玉律”]
这三个条件是任何声称求解 HJB 的数值方法都\emph{必须}验证的：单调性保证趋近正确的弱解（而非另一支）、稳定性防爆炸、一致性保证逼近正确方程。PINN/神经网络方法\emph{不}满足单调性——这正是它们没有收敛保证的根本原因。基于 DPP 直接离散的\textbf{半拉格朗日}格式 $V(x,t)\approx\min_a\{\Delta t\,L(x,a)+I_h[V](x+\Delta t\,f(x,a),t+\Delta t)\}$（$I_h$ 为网格插值）天然单调、且 CFL 无条件稳定（时间步不受空间网格约束），代价是精度依赖插值质量、高维插值复杂度指数增长。黏性解理论还提供\emph{稳定性}：若一列函数“几乎满足”HJB（残差小），则极限收敛到真黏性解（Barles--Perthame 半松弛极限）——这是数值近似可靠性的理论保障。
\end{remark}

\paragraph{静态 HJ 的快速算法与维数灾。}对 Eikonal 类 $|\nabla V|=f(x)$（$f>0$），利用其\emph{因果结构}（信息单向传播）：\textbf{Fast Marching}（Sethian 1996\cite{sethian1996fmm}）$\mathcal O(N\log N)$，是 Dijkstra 的连续版；\textbf{Fast Sweeping}（Zhao 2005\cite{zhao2005fsm}）$\mathcal O(N)$，交替方向扫描。机器人路径规划的 A\*/Dijkstra 正是 Fast Marching 的离散原型。但笛卡尔网格在 $n$ 维有 $\mathcal O(N^n)$ 点（\cref{tab:dp-num-curse}），$n\ge6$ 时即便有 GPU 集群也不可行。

\begin{table}[H]\centering\small
\caption{HJB 网格法的维数灾（每维 $N=100$）}
\label{tab:dp-num-curse}
\begin{tabular}{lccc}
\toprule
维度 $n$ & 网格点数 & 内存（float64） & 可行性 \\
\midrule
2 & $10^4$ & $80$\,KB & 毫秒 \\
3 & $10^6$ & $8$\,MB & 秒 \\
4 & $10^8$ & $800$\,MB & 分钟 \\
5 & $10^{10}$ & $80$\,GB & 需集群 \\
6 & $10^{12}$ & $8$\,TB & 极限 \\
$\ge7$ & $\ge10^{14}$ & --- & 不可行 \\
\bottomrule
\end{tabular}
\end{table}

\begin{remark}[E2 勘误：Hopf--Lax 的正确归属]
对 $\mathcal H$ 凸的情形，黏性解由经典的 \textbf{Hopf--Lax 公式}给出：$V(x,t)=\min_y\{J(y)+tH^*(\tfrac{x-y}{t})\}$（$H^*$ 为 $\mathcal H$ 的 Legendre 变换）。这一 min-plus 公式本身是 \textbf{Hopf 1965 / Lax} 的经典结果（$\mathcal H=\mathcal H(p)$ 凸、初值任意），\emph{不应}归于 Darbon--Osher 2016。Darbon--Osher 2016\cite{darbon2016hopf} 的真正贡献是用\emph{Hopf 公式}（对偶版，要求初值 $J$ 凸而 $\mathcal H$ 可\emph{非凸}）做逐点凸优化、把复杂度在维度上做到近线性，从而克服维数灾。两者措辞须分清。其余高维路线（稀疏网格、张量火车分解、DeepReach）仅列名不展开，详见 \cref{ch:plan-diffgame}。
\end{remark}

低维网格法的具体库（toolboxLS+helperOC、hj\_reachability、optimized\_dp、BEACLS 等）属工程实现，本章仅列名；博弈/可达性的数值实现（WENO/TVD-RK 的可达集实跑）见 \cref{sec:dg-numerics}，本节只给“单调格式通论 + Barles--Souganidis”的理论收敛保证。

% ════════════════════════════════════════════════════════════════
\section{与强化学习的逻辑接口}
\label{sec:dp-rl-bridge}

\paragraph{动机。}
本节用最短篇幅回答“DP 与 RL 是什么关系”，并把深入留给后续 RL 部。答案是：\textbf{RL 是 DP 在“无模型 + 函数近似 + 有限样本”约束下的近似求解方法}。

逻辑链 DP $\to$ TD $\to$ $Q$-learning 一目了然（\cref{tab:dp-rl-chain}）。关键是：TD(0) 更新 $V(s)\leftarrow V(s)+\alpha[r+\gamma V(s')-V(s)]$ 中的 $\delta=r+\gamma V(s')-V(s)$ 就是\textbf{Bellman 残差的单样本估计}，$\mathbb E[\delta\mid s]=(\mathcal TV)(s)-V(s)$。因此 TD(0) 是“随机梯度下降最小化 Bellman 残差”——几乎所有 value-based RL 算法都可理解为“最小化某种 Bellman 残差 $\mathbb E_{(s,a,r,s')}[(r+\gamma V_\theta(s')-V_\theta(s))^2]$”，DQN 用 target network 冻结右端（即 Banach 迭代的数值实现）、SAC 加熵正则、TD3 用 clipped double Q。

\begin{table}[H]\centering\small
\caption{从精确 DP 到深度 RL 的逻辑链（放松了哪些假设）}
\label{tab:dp-rl-chain}
\begin{tabular}{llll}
\toprule
层次 & 方法 & 知道什么 & Bellman 更新 \\
\midrule
精确 DP & VI / PI & 完整模型 $P,R$ & 精确计算 $\mathcal TV$ \\
Model-based RL & Dyna & 学得的 $\hat P,\hat R$ & 用估计模型做 VI \\
TD(0) & model-free & 无模型，只有样本 & $V\leftarrow V+\alpha[r+\gamma V'-V]$ \\
$Q$-learning & model-free & 无模型，只有样本 & $Q\leftarrow Q+\alpha[r+\gamma\max_{a'}Q'-Q]$ \\
DQN & NN-based & 无模型 + 函数近似 & $\min_\theta(r+\gamma\max_{a'}Q_{\theta^-}'-Q_\theta)^2$ \\
\bottomrule
\end{tabular}
\end{table}

连续时间也成立：critic 网络 $V_\theta(x)$ 的训练损失就是折扣 HJB 的残差 $\mathcal L_{\mathrm{critic}}=\tfrac12(\beta V_\theta-L-\nabla V_\theta\cdot f)^2$，最小化它即在求解 HJB——这是 PINN 用于值函数学习的理论基础（Doya 2000\cite{doya2000continuous} 的连续时间 RL / HJB 残差作为 TD 误差；注意 Doya 2000 是连续时间 TD/actor-critic + 值梯度贪心策略，论文并未以“$Q$-learning”命名，此即 E5 勘误）。

\begin{tcolorbox}[colback=second!4,colframe=second!55!black,title=接口提示：致命三角与深入去向]
当\textbf{函数近似 + bootstrapping + off-policy} 三者同时出现时，$\Pi\mathcal T$（投影 $\circ$ Bellman）可能不再是压缩、迭代可能发散（Sutton--Barto §11.3\cite{sutton2018rl}）——去掉任意一个都安全（tabular / Monte Carlo / on-policy）。这与 \cref{sec:dp-curse} 的 Baird/TvR 反例（E1）一脉相承。本节只搭接口，TD/$Q$-learning/Actor--Critic/离线 RL 的完整理论留给后续强化学习部展开。
\end{tcolorbox}

% ════════════════════════════════════════════════════════════════
\section{前沿一瞥}
\label{sec:dp-frontier}

本章止于单智能体确定性 HJB 与黏性解地基。三条延伸方向在此各点一句、给出去向，不写算法、不写代码。

\textbf{随机 HJB（Itô--Bellman）。}受控扩散 $dx_t=f\,dt+\sigma\,dW_t$ 的受控生成元含二阶项 $\mathcal L^uV=f\cdot\nabla V+\tfrac12\mathrm{tr}(\sigma\sigma^\top\nabla^2V)$，随机 HJB 为 $V_t+\min_u\{L+\mathcal L^uV\}=0$——二阶项正来自 Itô 公式的 $(dW)^2=dt$ 修正，须用二阶黏性解（semijets）。须注意：Stratonovich 积分下并非“DPP 失败”，而是无 $(dW)^2=dt$ 修正、需先转 Itô 形式（Wong--Zakai）后 DPP 才直接可用（E7 勘误）。完整随机 HJB 理论见 \cite{flemingsoner2006}。

\textbf{路径积分 / Kappen 线性化。}当 $\sigma\sigma^\top\propto \mathbf R^{-1}$（噪声协方差与控制代价成比例）时，随机 HJB 经对数变换\emph{精确线性化}，Feynman--Kac 给出路径积分解析解\cite{kappen2005linear,kappen2005path}——这是 MPPI 算法的理论支柱，但 MPPI 算法主体见 \cref{ch:mppi}。

\textbf{HJ 可达性 / 微分博弈。}面临对抗扰动时，单智能体 HJB 推广为 HJ-Isaacs 方程，可达集 = 值函数零水平集（DeepReach 用神经网络突破维数灾\cite{bansal2021deepreach}）。Isaacs 方程、微分博弈、上下值、后向可达管道 BRT、reach-avoid 与 $H_\infty$ 作为 min-max 博弈，均由 \cref{ch:plan-diffgame} 系统覆盖（BRT 方程归 Mitchell--Bayen--Tomlin 2005，作者含 Bayen，此即 E8 勘误\cite{mitchell2005time}），本章仅外引、不重写。

% ════════════════════════════════════════════════════════════════
\section*{本章小结}
\addcontentsline{toc}{section}{本章小结}

本章把最优控制的两条理论脊柱讲透：离散时间动态规划的\emph{一般理论}（压缩 / VI / PI / MPI / MDP / 维数灾难），与连续时间 HJB 的\emph{黏性解地基}。主线是“最优性原理 $\to$ Bellman 算子压缩 $\to$ 迭代算法 $\to$ 离散到连续极限 $\to$ HJB 与验证定理 $\to$ HJB$\leftrightarrow$PMP 对偶 $\to$ 经典解崩溃 $\to$ 黏性解”。其核心结论速查见 \cref{tab:dp-summary}。

\begin{table}[H]\centering\small
\caption{核心定理与方程速查}
\label{tab:dp-summary}
\begin{tabular}{lll}
\toprule
知识点 & 核心结论 & 出处 \\
\midrule
最优性原理 & 最优子策略最优（代价可加 + 马尔可夫） & \cref{thm:dp-optimality} \\
$\gamma$-压缩 & $\|\mathcal TV_1-\mathcal TV_2\|_\infty\le\gamma\|V_1-V_2\|_\infty$ & \cref{thm:dp-contraction} \\
VI 几何收敛 & $\|V_k-V^*\|_\infty\le\gamma^k\|V_0-V^*\|_\infty$ & \cref{cor:dp-banach} \\
PI 有限步终止 & $\le|\mathcal A|^{|\mathcal S|}$ 步；PI = Newton & \cref{thm:dp-pi-finite,thm:dp-ye-newton} \\
协态 = 梯度 & $\lambda_k=\nabla_xV_k(x_k^*)$ & \cref{thm:dp-costate} \\
HJB 方程 & $V_t+\mathcal H(x,\nabla V)=0$ & \cref{eq:dp-hjb} \\
验证定理 & 满足 HJB 的 $C^1$ 解即值函数（充分） & \cref{thm:dp-verification} \\
HJB$\leftrightarrow$PMP & $\lambda=\nabla V$；特征线 = Hamilton 方程 & \cref{thm:dp-duality} \\
黏性解 & 测试函数从上/下触碰；唯一性 & \cref{def:dp-viscosity,thm:dp-comparison} \\
值函数 = 唯一黏性解 & 求最优控制 $\equiv$ 求 HJB 黏性解 & \cref{thm:dp-perron} \\
单调格式收敛 & 单调 + 稳定 + 一致 $\Rightarrow$ 收敛 & \cref{thm:dp-barles} \\
\bottomrule
\end{tabular}
\end{table}

\section*{常见误解汇总}
\addcontentsline{toc}{section}{常见误解汇总}
\begin{itemize}
  \item \textbf{“最优性原理对所有最优策略的截断都成立。”}\;原理只保证\emph{存在}一个最优截断；并列冗余的最优策略不受约束。
  \item \textbf{“$\gamma=1$ 直接套压缩性。”}\;$\gamma=1$ 时算子仅非扩张，不保证唯一不动点；须 proper 策略或加权范数。
  \item \textbf{“NN 自动克服维数灾。”}\;统计样本复杂度仍随 $d$ 增长；线性近似下 $Q$-learning 发散反例是 Baird 1995、off-policy TD 发散是 TvR 1997，二者不可混（E1）。
  \item \textbf{“HJB 总有 $C^1$ 解。”}\;一般只有 Lipschitz 黏性解；用标准 FEM/Galerkin 会震荡，须 upwind/单调格式（E 见 \cref{thm:dp-barles}）。
  \item \textbf{“Hopf--Lax 公式归 Darbon--Osher 2016。”}\;该 min-plus 公式是 Hopf 1965/Lax 经典结果；Darbon--Osher 用的是 Hopf 公式做高维凸优化（E2）。
  \item \textbf{“Stratonovich 下 DPP 失败。”}\;只是无 $(dW)^2=dt$ 修正、须先转 Itô，转换后 DPP 仍可用（E7）。
\end{itemize}

\section*{延伸阅读}
\addcontentsline{toc}{section}{延伸阅读}
离散 DP 与 MDP 的权威是 Bertsekas \emph{DP \& Optimal Control}\cite{bertsekas2017dpoc1} 与 Puterman \emph{MDP}\cite{puterman1994mdp}；RL 视角下的 DP 见 Sutton--Barto\cite{sutton2018rl}；近似 DP 的误差理论见 Munos--Szepesvári\cite{munos2008finite} 与 Powell\cite{powell2011adp}。连续 HJB 与黏性解的标准教材是 Bardi--Capuzzo-Dolcetta\cite{bardi1997optimal} 与 Fleming--Soner\cite{flemingsoner2006}，比较原理与二阶 semijets 的“用户指南”是 Crandall--Ishii--Lions\cite{crandall1992users}，数值收敛地基是 Barles--Souganidis\cite{barles1991convergence}。机器人语境的 DP/HJB 应用见 Tedrake \emph{Underactuated Robotics}\cite{tedrake_underactuated}。

\section*{本章与后续章节的关系}
\addcontentsline{toc}{section}{本章与后续章节的关系}
本章是卷四《控制理论基础》最优控制脊柱的“场级地基”。承上：它补全了 \cref{ch:control,ch:optimal-control-basics} 留下的 DP 一般理论与黏性解两块空白，并依赖 \cref{ch:opt-variational-pmp}（变分/PMP）的 Hamiltonian 与必要条件（二章在 Hamiltonian 符号约定上双向接口）。启下：\cref{ch:lqr-lqg-riccati}（LQR/LQG/Riccati 深化）接手完整 Riccati、确定性等价与时域 $H_\infty$；\cref{ch:lyapunov-stability}（Lyapunov/ISS）承接“HJB 不等式 = Lyapunov 下界”；\cref{ch:clf-cbf}（CLF/CBF）把“HJB 极小化条件的单点投影”发展为安全控制；\cref{ch:ddp-ilqr}（DDP/iLQR）把“局部二次 DP”工程化；\cref{ch:mpc-advanced}（高级 MPC）把退行视界 DP 与终端值函数做深；而 \cref{ch:hj-reachability} 与 \cref{ch:plan-diffgame} 把单智能体 HJB 推广到可达性与微分博弈。沿着本章铺好的“值函数 = 唯一黏性解”这一地基，最优、鲁棒与安全控制在卷四后续诸章次第展开。
